\frfilename{mt252.tex}
\versiondate{6.12.07}
\copyrightdate{2000}

\def\displaycause#1{\noalign{\noindent (#1)}}

\def\chaptername{Ukuran produk}
\def\sectionname{Teorema Fubini}

\newsection{252}

Barangkali ciri terpenting dari konsep `ukuran produk' adalah bahwa kita
dapat menggunakannya untuk membahas integral berulang. Dalam bagian ini
saya memberikan versi teorema Fubini dan teorema Tonelli (252B, 252G)
beserta berbagai akibatnya; yang paling berguna adalah versi untuk ruang
$\sigma$-hingga (252C, 252H). Sebagai penerapan, saya menguraikan hubungan
antara integrasi dan pengukuran himpunan ordinat (252N), serta menghitung
volume berdimensi $r$ dari sebuah bola di $\BbbR^r$
(252Q\cmmnt{, 252Xi}).
Saya menyebutkan contoh-contoh tandingan yang menunjukkan kesulitan yang
dapat timbul pada ukuran yang tidak $\sigma$-hingga dan fungsi yang tidak
terintegralkan (252K-252L\cmmnt{, 252Xf-252Xg}).

\leader{252A}{Integral berulang} Misalkan $(X,\Sigma,\mu)$ dan
$(Y,\Tau,\nu)$ adalah ruang ukur, dan $f$ suatu fungsi bernilai real yang
didefinisikan pada suatu himpunan $\dom f\subseteq X\times Y$. Kita dapat
berusaha membentuk {\bf integral berulang}

\Centerline{$\iint f(x,y)\nu(dy)\mu(dx)
=\int\bigl(\int f(x,y)\nu(dy)\bigr)\mu(dx)$,}

\noindent yang harus ditafsirkan sebagai berikut: tetapkan

\Centerline{$D=\{x:x\in X,\,\int f(x,y)\nu(dy)$ terdefinisi dalam
$[-\infty,\infty]\}$,}

\Centerline{$g(x)=\int f(x,y)\nu(dy)$ untuk $x\in D$,}

\noindent lalu tuliskan $\iint f(x,y)\nu(dy)\mu(dx)=\int g(x)\mu(dx)$
jika integral terakhir ini terdefinisi. \cmmnt{Tentu saja, subhimpunan
$Y$ tempat $y\mapsto f(x,y)$ didefinisikan dapat berubah menurut $x$,
tetapi subhimpunan itu harus selalu koterabaikan, demikian pula $D$.

Dengan cara yang sama, dengan menukar peran $X$ dan $Y$, kita dapat
berusaha membentuk integral berulang

\Centerline{$\iint f(x,y)\mu(dx)\nu(dy)
=\int\bigl(\int f(x,y)\mu(dx)\bigr)\nu(dy)$.}

\noindent Intinya ialah bahwa, dengan syarat yang sesuai pada $\mu$ dan
$\nu$, kita dapat menghubungkan kedua integral berulang ini satu sama lain
dengan menghubungkan keduanya dengan integral $f$ sendiri terhadap ukuran
produk pada $X\times Y$.

Sebagaimana akan segera tampak, di sini sangat penting untuk mengizinkan
pembahasan integral suatu fungsi yang tidak didefinisikan di setiap titik.
Tidak terlalu penting apakah integran dan integral diizinkan bernilai tak
hingga, tetapi agar tegas, saya akan mengikuti aturan 135F; yaitu,
$\int f=\int f^+-\int f^-$ asalkan $f$ didefinisikan hampir di mana-mana,
bernilai dalam $[-\infty,\infty]$, terukur secara virtual, dan paling banyak
salah satu dari $\int f^+$, $\int f^-$ bernilai tak hingga.
}%end of comment

\leader{252B}{Teorema} Misalkan $(X,\Sigma,\mu)$ dan $(Y,\Tau,\nu)$
adalah ruang ukur, dengan produk c.l.d.\ $(X\times Y,\Lambda,\lambda)$
\cmmnt{ (251F)}. Andaikan $\nu$
$\sigma$-hingga dan $\mu$ {\it dapat dilokalkan secara ketat} atau
{\it lengkap dan ditentukan secara lokal}. Misalkan $f$ suatu fungsi
bernilai $[-\infty,\infty]$ sedemikian sehingga $\int fd\lambda$
terdefinisi dalam $[-\infty,\infty]$. Maka
$\iint f(x,y)\nu(dy)\mu(dx)$ terdefinisi dan sama dengan
$\int fd\lambda$.

\proof{ Pembuktian hasil ini melibatkan kesulitan teknis yang cukup besar.
Jika Anda belum pernah menjumpai gagasan-gagasan ini, hampir pasti Anda
tidak sebaiknya langsung menghadapi seluruh keluasan versi yang dinyatakan
di atas. Karena itu, saya akan memulai dengan menuliskan pembuktian untuk
kasus ketika $\mu$ dan $\nu$ keduanya hingga total; kasus ini saja sudah
cukup panjang. Saya akan menyajikannya sedemikian rupa sehingga hanya
bagian tengah (bagian (b) di bawah) yang perlu diubah dalam kasus umum.
Kemudian, setelah menyelesaikan pembuktian kasus khusus, saya akan
memberikan versi alternatif dari (b) yang diperlukan untuk hasil lengkap.

\medskip

{\bf (a)} Tuliskan $\eusm L$ untuk keluarga semua fungsi bernilai
$[0,\infty]$ yang, bila dilambangkan dengan $f$, membuat $\int fd\lambda$ dan
\discrversionA{\break}{}$\iint f(x,y)\nu(dy)\mu(dx)$ terdefinisi dan
sama. Tujuan saya mula-mula ialah menunjukkan bahwa $f\in\eusm L$ setiap
kali $f$ nonnegatif dan $\int fd\lambda$ terdefinisi, lalu menelaah selisih
fungsi-fungsi dalam $\eusm L$. Untuk membuktikan bahwa cukup banyak fungsi
termasuk dalam $\eusm L$, strategi saya ialah memulai dengan fungsi
`elementer' dan bergerak keluar melalui kelas-kelas yang makin besar.
Cara yang paling efisien ialah memulai dengan menguraikan cara membangun
anggota baru $\eusm L$ dari anggota lama, sebagai berikut.

\medskip

\quad{\bf (i)} $f_1+f_2\in\eusm L$ untuk semua $f_1$, $f_2\in\eusm L$,
dan $cf\in\eusm L$ untuk semua $f\in\eusm L$,
$c\in\coint{0,\infty}$; ini karena

\Centerline{$\int(f_1+f_2)(x,y)\nu(dy)
=\int f_1(x,y)\nu(dy)+\int f_2(x,y)\nu(dy)$,}

\Centerline{$\int (cf)(x,y)\nu(dy)=c\int f(x,y)\nu(dy)$}

\noindent setiap kali ruas-ruas kanan terdefinisi, yang menurut andaian
berlaku untuk hampir setiap $x$, sehingga

$$\eqalign{\iint(f_1+f_2)(x,y)\nu(dy)\mu(dx)
&=\iint f_1(x,y)\nu(dy)\mu(dx)+\iint f_2(x,y)\nu(dy)\mu(dx)\cr
&=\int f_1d\lambda+\int f_2d\lambda
=\int(f_1+f_2)d\lambda,\cr}$$

$$\iint(cf)(x,y)\nu(dy)\mu(dx)
=c\int f(x,y)\nu(dy)\mu(dx)
=c\int fd\lambda
=\int(cf)d\lambda.$$

\medskip

\quad{\bf (ii)} Jika $\sequencen{f_n}$ adalah barisan dalam $\eusm L$
sedemikian sehingga $f_n(x,y)\le f_{n+1}(x,y)$ setiap kali
$n\in\Bbb N$ dan $(x,y)\in\dom f_n\cap\dom f_{n+1}$, maka
$\sup_{n\in\Bbb N}f_n\in\eusm L$. \Prf\
Tetapkan $f=\sup_{n\in\Bbb N}f_n$; untuk $x\in X$, $n\in\Bbb N$,
tetapkan $g_n(x)=\int f_n(x,y)\nu(dy)$ apabila integral ini terdefinisi
dalam $[0,\infty]$. Karena di sini saya mengizinkan $\infty$ sebagai nilai
suatu fungsi, wajar untuk menganggap $f$ didefinisikan pada
$\bigcap_{n\in\Bbb N}\dom f_n$. Menurut teorema B.Levi,
$\int fd\lambda=\sup_{n\in\Bbb N}\int f_nd\lambda$; tuliskan $u$ untuk
nilai bersama ini dalam $[0,\infty]$. Selanjutnya, karena
$f_n\le f_{n+1}$ di setiap titik tempat keduanya didefinisikan,
$g_n\le g_{n+1}$ di setiap titik tempat keduanya didefinisikan, untuk
setiap $n$; menurut andaian $f_n\in\eusm L$, jadi $g_n$ didefinisikan
$\mu$-hampir di mana-mana untuk setiap $n$, dan

\Centerline{$\sup_{n\in\Bbb N}\int g_nd\mu
=\sup_{n\in\Bbb N}\int f_nd\lambda=u$.}

\noindent Dengan menerapkan teorema B.Levi lagi, $\int g\,d\mu=u$, dengan
$g=\sup_{n\in\Bbb N}g_n$. Sekarang ambil sembarang
$x\in\bigcap_{n\in\Bbb N}\dom g_n$, dan tinjau fungsi-fungsi $f_{xn}$
pada $Y$, dengan menetapkan $f_{xn}(y)=f_n(x,y)$ setiap kali ruas kanan
ini didefinisikan. Setiap $f_{xn}$ mempunyai integral dalam
$[0,\infty]$, dan $f_{xn}(y)\le f_{x,n+1}(y)$ setiap kali keduanya
didefinisikan, serta

\Centerline{$\sup_{n\in\Bbb N}\int f_{xn}d\nu=g(x)$;}

\noindent maka, dengan menggunakan teorema B.Levi untuk ketiga kalinya,
$\int(\sup_{n\in\Bbb N}f_{xn})d\nu$ terdefinisi dan sama dengan $g(x)$,
yakni

\Centerline{$\int f(x,y)\nu(dy)=g(x)$.}

\noindent Hal ini berlaku untuk hampir setiap $x$, sehingga

\Centerline{$\iint f(x,y)\nu(dy)\mu(dx)
=\int g\,d\mu=u=\int fd\lambda$.}

\noindent Jadi $f\in\eusm L$, sebagaimana dinyatakan.\ \Qed

\medskip

\quad{\bf (iii)} Pengungkapan gagasan-gagasan pada bagian pembuktian
berikutnya akan lebih lancar jika saya memperkenalkan istilah lain.
Tuliskan $\Cal W$ untuk
$\{W:W\subseteq X\times Y,\,\chi W\in\eusm L\}$. Maka

\inset{($\alpha$) jika $W$, $W'\in\Cal W$ dan $W\cap W'=\emptyset$,
$W\cup W'\in\Cal W$}

\noindent menurut (i), karena $\chi(W\cup W')=\chi W+\chi W'$,

\inset{($\beta$) $\bigcup_{n\in\Bbb N}W_n\in\Cal W$ setiap kali
$\sequencen{W_n}$ merupakan barisan tak-menurun dalam $\Cal W$}

\noindent karena $\sequencen{\chi W_n}\uparrow\chi W$, dan kita dapat
menggunakan (ii).

Perlu dicatat pula bahwa, untuk setiap $W\subseteq X\times Y$ dan setiap
$x\in X$, $\int\chi W(x,y)\nu(dy)=\nu W[\{x\}]$, setidaknya apabila
$W[\{x\}]=\{y:(x,y)\in W\}$ terukur oleh $\nu$. Selain itu, karena
$\lambda$ lengkap, suatu himpunan $W\subseteq X\times Y$ termasuk dalam
$\Lambda$ jika dan hanya jika $\chi W$ terukur secara virtual terhadap
$\lambda$, jika dan hanya jika $\int\chi W\,d\lambda$ terdefinisi dalam
$[0,\infty]$; dan dalam hal ini
$\lambda W=\int\chi W\,d\lambda$.

\medskip

\quad{\bf (iv)} Akhirnya, kita perlu mengamati bahwa, dalam keadaan yang
sesuai, selisih dua anggota $\Cal W$ akan termasuk dalam $\Cal W$: jika
$W$, $W'\in\Cal W$, $W\subseteq W'$, dan $\lambda W'<\infty$, maka
$W'\setminus W\in\Cal W$. \Prf\ Menurut andaian,
$g(x)=\int\chi W(x,y)\nu(dy)$ dan
$g'(x)=\int\chi W'(x,y)\nu(dy)$ terdefinisi untuk hampir setiap $x$,
serta $\int g\,d\mu=\lambda W$, $\int g'd\mu=\lambda W'$. Karena
$\lambda W'$ berhingga, $g'$ harus berhingga hampir di mana-mana, dan
$D=\{x:x\in\dom g\cap\dom g',\,g'(x)<\infty\}$ koterabaikan. Sekarang,
untuk setiap $x\in D$, $g(x)$ dan $g'(x)$ keduanya berhingga, sehingga

\Centerline{$y\mapsto\chi(W'\setminus W)(x,y)=\chi W'(x,y)-\chi W(x,y)$}

\noindent adalah selisih dua fungsi terintegralkan, dan

$$\eqalign{\int\chi(W'\setminus W)(x,y)\nu(dy)
&=\int\chi W'(x,y)-\chi W(x,y)\nu(dy)\cr
&=\int\chi W'(x,y)\nu(dy)-\int\chi W(x,y)\nu(dy)
=g'(x)-g(x).\cr}$$

\noindent Dengan demikian

\Centerline{$\iint\chi(W'\setminus W)(x,y)\nu(dy)\mu(dx)
=\int g'(x)-g(x)\mu(dx)
=\lambda W'-\lambda W
=\lambda(W'\setminus W)$,}

\noindent dan $W'\setminus W$ termasuk dalam $\Cal W$.\ \Qed

(Tentu saja, argumen tepat di atas dapat dipersingkat beberapa kata jika
kita mengizinkan diri untuk mengandaikan bahwa $\mu$ dan $\nu$ hingga
total, sebab dalam hal itu $g(x)$ dan $g'(x)$ akan berhingga setiap kali
keduanya didefinisikan; tetapi gagasan kuncinya, bahwa selisih fungsi-fungsi
terintegralkan adalah terintegralkan, tidak berubah.)

\medskip

{\bf (b)} Sekarang mari kita telaah kelas $\Cal W$, dengan mengandaikan
bahwa $\mu$ dan $\nu$ hingga total.

\medskip

\quad{\bf (i)} $E\times F\in\Cal W$ untuk semua $E\in\Sigma$,
$F\in\Tau$. \Prf\ $\lambda(E\times F)=\mu E\cdot\nu F$ (251J), dan

\Centerline{$\int\chi(E\times F)(x,y)\nu(dy)=\nu F\,\chi E(x)$}

\noindent untuk setiap $x$, sehingga

$$\eqalign{\iint\chi(E\times F)(x,y)\nu(dy)\mu(dx)
&=\int(\nu F\,\chi E(x))\mu(dx)
=\mu E\cdot\nu F\cr
&=\lambda(E\times F)
=\int\chi(E\times F)d\lambda. \text{ \Qed}\cr}$$

\medskip

\quad{\bf (ii)} Misalkan $\Cal C$ adalah
$\{E\times F:E\in\Sigma,\,F\in\Tau\}$.
Maka $\Cal C$ tertutup terhadap irisan berhingga (karena
$(E\times F)\cap(E'\times F')=(E\cap E')\times(F\cap F')$)
dan termuat dalam $\Cal W$. Khususnya, $X\times Y\in\Cal W$. Akan tetapi,
hal ini, bersama dengan (a-iv) dan (a-iii-$\beta$) di atas, berarti bahwa
$\Cal W$ merupakan kelas Dynkin (definisi: 136A), sehingga memuat
aljabar-sigma subhimpunan $X\times Y$ yang dibangkitkan oleh $\Cal C$,
menurut Teorema Kelas Monoton (136B); yaitu,
$\Cal W\supseteq\Sigma\tensorhat\Tau$ (definisi: 251D).

\medskip

\quad{\bf (iii)} Selanjutnya, $W\in\Cal W$ setiap kali
$W\subseteq X\times Y$ terabaikan terhadap $\lambda$. \Prf\ Menurut
251Ib, terdapat $V\in\Sigma\tensorhat\Tau$ sedemikian sehingga
$V\subseteq(X\times Y)\setminus W$ dan
$\lambda V=\lambda((X\times Y)\setminus W)$. Karena
$\lambda(X\times Y)=\mu X\cdot\nu Y$ berhingga,
$V'=(X\times Y)\setminus V$ terabaikan terhadap $\lambda$, dan kita
mempunyai $W\subseteq V'\in\Sigma\tensorhat\Tau$. Akibatnya

\Centerline{$0=\lambda V'=\iint\chi V'(x,y)\nu(dy)\mu(dx)$.}

\noindent Akan tetapi, ini berarti bahwa

\Centerline{$D=\{x:\int\chi V'(x,y)\nu(dy)$ terdefinisi dan sama
dengan $0\}$}

\noindent koterabaikan. Jika $x\in D$, kita harus mempunyai
$\chi V'(x,y)=0$ untuk $\nu$-hampir setiap $y$, yaitu $V'[\{x\}]$
terabaikan; dalam hal ini $W[\{x\}]\subseteq V'[\{x\}]$ juga terabaikan,
dan $\int\chi W(x,y)\nu(dy)=0$. Hal ini berlaku untuk setiap $x\in D$,
jadi $\int\chi W(x,y)\nu(dy)$ terdefinisi dan sama dengan $0$ untuk
hampir setiap $x$, dan

\Centerline{$\iint\chi W(x,y)\nu(dy)\mu(dx)=0=\lambda W$,}

\noindent sebagaimana diperlukan.\ \Qed

\medskip

\quad{\bf (iv)} Maka $\Lambda\subseteq\Cal W$. \Prf\ Jika
$W\in\Lambda$, maka, sekali lagi menurut 251Ib, terdapat
$V\in\Sigma\tensorhat\Tau$ sedemikian sehingga $V\subseteq W$ dan
$\lambda V=\lambda W$, sehingga $\lambda(W\setminus V)=0$. Sekarang
$V\in\Cal W$ menurut (ii) dan $W\setminus V\in\Cal W$ menurut (iii),
jadi $W\in\Cal W$ menurut (a-iii-$\alpha$).\ \Qed

\medskip

{\bf (c)} Saya kembali ke kelas $\eusm L$.

\medskip

\quad{\bf (i)} Jika $f\in\eusm L$ dan $g$ adalah fungsi bernilai
$[0,\infty]$ yang didefinisikan dan sama dengan $f\,\,\lambda$-hampir
di mana-mana, maka $g\in\eusm L$. \Prf\ Tetapkan

\Centerline{$W=(X\times Y)\setminus\{(x,y):(x,y)\in\dom f\cap\dom g$,
$f(x,y)=g(x,y)\}$,}

\noindent sehingga $\lambda W=0$. (Ingat bahwa $\lambda$ lengkap.)
Menurut (b), $\iint\chi W(x,y)\nu(dy)\mu(dx)=0$, yaitu $W[\{x\}]$
terabaikan terhadap $\nu$ untuk $\mu$-hampir setiap $x$. Misalkan $D$
adalah $\{x:x\in X,\,W[\{x\}]$ terabaikan terhadap $\nu$ itu$\}$. Maka $D$
koterabaikan terhadap $\mu$. Jika $x\in D$, maka

\Centerline{$W[\{x\}]=Y\setminus\{y:(x,y)\in\dom f\cap\dom g$,
$f(x,y)=g(x,y)\}$}

\noindent terabaikan, sehingga
$\int f(x,y)\nu(dy)=\int g(x,y)\nu(dy)$ jika salah satunya terdefinisi.
Jadi fungsi-fungsi
\discrcenter{468pt}{$x\mapsto\int f(x,y)\nu(dy)$,
\ifdim\pagewidth>467pt\quad\fi
$x\mapsto\int g(x,y)\nu(dy)$ }sama hampir di mana-mana, dan

\Centerline{$\iint g(x,y)\nu(dy)\mu(dx)=\iint
f(x,y)\nu(dy)\mu(dx)
=\int fd\lambda=\int g\,d\lambda$,}

\noindent sehingga $g\in\eusm L$.\ \Qed

\medskip

\quad{\bf (ii)} Sekarang misalkan $f$ sembarang fungsi nonnegatif
sedemikian sehingga $\int fd\lambda$ terdefinisi dalam $[0,\infty]$.
Maka $f\in\eusm L$. \Prf\ Untuk $k$, $n\in\Bbb N$, tetapkan

\Centerline{$W_{nk}=\{(x,y):(x,y)\in\dom f,\,f(x,y)\ge 2^{-n}k\}$.}

\noindent Karena $\lambda$ lengkap, $f$ terukur secara virtual terhadap
$\lambda$, dan $\dom f$ koterabaikan, setiap $W_{nk}$ termasuk dalam
$\Lambda$, jadi $\chi W_{nk}\in\eusm L$ menurut (b). Tetapkan
$f_n=\sum_{k=1}^{4^n}2^{-n}\chi W_{nk}$, sehingga

$$\eqalign{f_n(x,y)
&=2^{-n}k\text{ jika }k\le 4^n\text{ dan }2^{-n}k
                   \le f(x,y)<2^{-n}(k+1),\cr
&=2^n\text{ jika }f(x,y)\ge 2^n,\cr
&=0\text{ jika }(x,y)\in(X\times Y)\setminus\dom f.\cr}$$

\noindent Menurut (a-i), $f_n\in\eusm L$ untuk setiap $n\in\Bbb N$,
sedangkan $\sequencen{f_n}$ tak-menurun, jadi
$f'=\sup_{n\in\Bbb N}f_n\in\eusm L$ menurut (a-ii). Namun
$f\eae f'$, jadi $f\in\eusm L$ menurut (i) tepat di atas.\
\Qed

\medskip

\quad{\bf (iii)} Akhirnya, misalkan $f$ sembarang fungsi bernilai
$[-\infty,\infty]$ sedemikian sehingga $\int fd\lambda$ terdefinisi
dalam $[-\infty,\infty]$. Maka $\int f^+d\lambda$ dan
$\int f^-d\lambda$ keduanya terdefinisi dan paling banyak salah satunya
bernilai tak hingga. Menurut (ii), $f^+$ dan $f^-$ keduanya termasuk dalam
$\eusm L$. Tetapkan $g(x)=\int f^+(x,y)\nu(dy)$,
$h(x)=\int f^-(x,y)\nu(dy)$ setiap kali ruas-ruas ini terdefinisi; maka
$\int g\,d\mu=\int f^+d\lambda$ dan $\int h\,d\mu=\int f^-d\lambda$
keduanya terdefinisi dalam $[0,\infty]$.

Andaikan terlebih dahulu bahwa $\int f^-d\lambda$ berhingga. Maka
$\int h\,d\mu$ berhingga, sehingga $h$ harus berhingga
$\mu$-hampir di mana-mana; tetapkan

\Centerline{$D=\{x:x\in\dom g\cap\dom h,\,h(x)<\infty\}$.}

\noindent Untuk setiap $x\in D$, $\int f^+(x,y)\nu(dy)$ dan
$\int f^-(x,y)\nu(dy)$ terdefinisi dalam $[0,\infty]$, dan yang terakhir
berhingga; jadi

\Centerline{$\int f(x,y)\nu(dy)
=\int f^+(x,y)\nu(dy)-\int f^-(x,y)\nu(dy)
=g(x)-h(x)$}

\noindent terdefinisi dalam $\ocint{-\infty,\infty}$. Karena $D$
koterabaikan,

$$\eqalign{\iint f(x,y)\nu(dy)\mu(dx)
&=\int g(x)-h(x)\mu(dx)
=\int g\,d\mu-\int h\,d\mu\cr
&=\int f^+d\lambda-\int f^-d\lambda
=\int fd\lambda,\cr}$$

\noindent sebagaimana diperlukan.

Jadi kita telah memperoleh hasilnya apabila $\int f^-d\lambda$ berhingga.
Dengan cara yang sama, atau dengan menerapkan argumen di atas pada $-f$,
kita melihat bahwa $\iint f(x,y)\nu(dy)\mu(dx)=\int fd\lambda$
apabila $\int f^+d\lambda$ berhingga.

Dengan demikian teorema ini telah terbukti, setidaknya ketika $\mu$ dan
$\nu$ hingga total.

\medskip

{\bf (b*)} Satu-satunya titik dalam argumen di atas tempat kita perlu
mengetahui sesuatu yang khusus tentang ukuran $\mu$ dan $\nu$ adalah
bagian (b), ketika menunjukkan bahwa $\Lambda\subseteq\Cal W$. Sekarang
saya kembali ke titik ini di bawah hipotesis teorema sebagaimana dinyatakan,
yaitu bahwa $\nu$ $\sigma$-hingga dan $\mu$ dapat dilokalkan secara ketat,
atau lengkap dan ditentukan secara lokal.

\medskip

\quad{\bf (i)} Perlu dicatat bahwa pelengkapan $\hat\mu$ dari $\mu$
(212C) identik dengan versi c.l.d.\ dari ukuran tersebut, yaitu
$\tilde\mu$ (213E). \Prf\ Jika
$\mu$ dapat dilokalkan secara ketat, maka $\hat\mu=\tilde\mu$ menurut
213Ha. Jika $\mu$ lengkap dan ditentukan secara lokal, maka
$\hat\mu=\mu=\tilde\mu$ (212D, 213Hf).\ \Qed

\medskip

\quad{\bf (ii)} Tuliskan
$\Sigma^f=\{G:G\in\Sigma,\,\mu G<\infty\}$,
$\Tau^f=\{H:H\in\Tau,\,\nu H<\infty\}$. Untuk $G\in\Sigma^f$,
$H\in\Tau^f$, misalkan $\mu_G$, $\nu_H$, dan
$\lambda_{G\times H}$ berturut-turut adalah ukuran subruang pada $G$,
$H$, dan $G\times H$; maka $\lambda_{G\times H}$ adalah ukuran produk
c.l.d.\ dari $\mu_G$ dan $\nu_H$ (251Q(ii-$\alpha$)). Sekarang
$W\cap(G\times H)\in\Cal W$ untuk setiap $W\in\Lambda$. \Prf\ $W\cap(G\times H)$
termasuk dalam domain $\lambda_{G\times H}$, jadi
menurut (b) dalam pembuktian ini, yang diterapkan pada ukuran hingga total
$\mu_G$ dan $\nu_H$,

$$\eqalignno{\lambda(W\cap(G\times H))
&=\lambda_{G\times H}(W\cap(G\times H))\cr
&=\int_G\int_H\chi(W\cap(G\times H))(x,y)\nu_H(dy)\mu_G(dx)\cr
&=\int_G\int_Y\chi(W\cap(G\times H))(x,y)\nu(dy)\mu_G(dx)\cr
\displaycause{karena $\chi(W\cap(G\times H))(x,y)=0$ jika
$y\in Y\setminus H$, sehingga kita dapat menggunakan 131E}
&=\int_X\int_Y\chi(W\cap(G\times H))(x,y)\nu(dy)\mu(dx)\cr}$$

\noindent menurut 131E lagi, karena
$\int_Y\chi(W\cap(G\times H))(x,y)\nu(dy)=0$ jika
$x\in X\setminus G$. Jadi $W\cap(G\times H)\in\Cal W$.\ \Qed

\medskip

\quad{\bf (iii)} Sesungguhnya, $W\in\Cal W$ untuk setiap
$W\in\Lambda$. \Prf\
Ingat bahwa kita mengandaikan $\nu$
$\sigma$-hingga. Misalkan $\sequencen{Y_n}$ adalah barisan tak-menurun
dalam $\Tau^f$ yang menutupi $Y$, dan untuk setiap $n\in\Bbb N$, tetapkan
$W_n=W\cap(X\times Y_n)$,
$g_n(x)=\int\chi W_n(x,y)\nu(dy)$ setiap kali integral ini terdefinisi.
Untuk setiap $G\in\Sigma^f$,

\Centerline{$\int_Gg_nd\mu
=\iint\chi(W\cap(G\times Y_n))(x,y)\nu(dy)\mu(dx)$}

\noindent terdefinisi dan sama dengan $\lambda(W\cap(G\times Y_n))$
menurut (ii). Namun ini berarti, pertama, bahwa
$G\setminus\dom g_n$ terabaikan, yaitu
$\hat\mu(G\setminus\dom g_n)=0$. Karena hal ini berlaku setiap kali
$\mu G$ berhingga, $\tilde\mu(X\setminus\dom g_n)=0$, dan $g_n$
didefinisikan $\tilde\mu$-hampir di mana-mana; tetapi
$\tilde\mu=\hat\mu$, sehingga $g_n$ didefinisikan
$\hat\mu$-hampir di mana-mana, yaitu $\mu$-hampir di mana-mana\ (212Eb).
Selanjutnya, jika kita menetapkan
$E_{na}=\{x:x\in\dom g_n,\,g_n(x)\ge a\}$ untuk $a\in\Bbb R$, maka
$E_{na}\cap G\in\hat\Sigma$ setiap kali $G\in\Sigma^f$, dengan
$\hat\Sigma$ domain $\hat\mu$; menurut definisi dalam 213D, $E_{na}$
terukur oleh $\tilde\mu=\hat\mu$. Karena $a$ sebarang, $g_n$ terukur
secara virtual terhadap $\mu$ (212Fa).

Dengan demikian, kita dapat membicarakan $\int g_nd\mu$. Sekarang

$$\eqalignno{\iint\chi W_n(x,y)\nu(dy)\mu(dx)
&=\int g_nd\mu
=\sup_{G\in\Sigma^f}\int_Gg_n\cr
\displaycause{213B, karena $\mu$ semihingga}
&=\sup_{G\in\Sigma^f}\lambda(W\cap(G\times Y_n))
=\lambda(W\cap(X\times Y_n))\cr}$$

\noindent menurut definisi dalam 251F. Jadi
$W\cap(X\times Y_n)\in\Cal W$.

Hal ini berlaku untuk setiap $n\in\Bbb N$. Karena
$\sequencen{Y_n}\uparrow Y$, $W\in\Cal W$ menurut (a-iii-$\beta$).\ \Qed

\medskip

\quad{\bf (iv)} Karena itu, kita dapat kembali ke bagian (c) dari
argumen di atas dan menarik kesimpulan seperti sebelumnya.
}%end of proof of 252B

\leader{252C}{}\cmmnt{ Teorema di atas tentu saja asimetris, dalam arti
bahwa hipotesis yang berbeda dikenakan pada kedua ukuran faktor $\mu$ dan
$\nu$. Jika kita menginginkan teorema `simetris', kita harus mengandaikan
bahwa keduanya $\sigma$-hingga, sebagai berikut.

\medskip

\noindent}{\bf Akibat} Misalkan $(X,\Sigma,\mu)$ dan $(Y,\Tau,\nu)$
adalah dua ruang ukur $\sigma$-hingga, dan $\lambda$ ukuran produk
c.l.d.\ pada $X\times Y$. Jika $f$ terintegralkan terhadap $\lambda$, maka
$\iint f(x,y)\nu(dy)\mu(dx)$ dan $\iint f(x,y)\mu(dx)\nu(dy)$
terdefinisi, berhingga, dan sama.

\proof{ Karena $\mu$ dan $\nu$ pasti dapat dilokalkan secara ketat
(211Lc), kita dapat menerapkan 252B dari kedua sisi untuk menyimpulkan bahwa

\Centerline{$\iint f(x,y)\nu(dy)\mu(dx)
=\int f\,d\lambda
=\iint f(x,y)\mu(dx)\nu(dy)$.}
}%end of proof of 252C

\leader{252D}{}\cmmnt{ Begitu banyak penerapan teorema Fubini menggunakan
fungsi indikator, sehingga saya mengambil beberapa baris untuk menjabarkan
bentuk 252B dalam kasus ini, seperti pada bagian (b)-(b*) pembuktiannya.

\medskip

\noindent}{\bf Akibat} Misalkan $(X,\Sigma,\mu)$ dan $(Y,\Tau,\nu)$
adalah ruang ukur dan $\lambda$ ukuran produk c.l.d.\ pada $X\times Y$.
Andaikan $\nu$ $\sigma$-hingga dan $\mu$ dapat dilokalkan secara ketat,
atau lengkap dan ditentukan secara lokal.

(i) Jika $W\in\dom\lambda$, maka $\int\nu^*W[\{x\}]\mu(dx)$
terdefinisi dalam $[0,\infty]$ dan sama dengan $\lambda W$.

(ii) Jika $\nu$ lengkap, kita dapat menuliskan
$\int\nu W[\{x\}]\mu(dx)$ sebagai pengganti
$\int\nu^*W[\{x\}]\mu(dx)$.

\proof{ Intinya hanyalah bahwa
$\int\chi W(x,y)\nu(dy)=\hat\nu W[\{x\}]$ setiap kali salah satu ruas
terdefinisi, dengan $\hat\nu$ pelengkapan $\nu$ (212F). Sekarang 252B
memberi tahu kita bahwa

\Centerline{$\lambda W=\iint\chi W(x,y)\nu(dy)\mu(dx)
=\int\hat\nu W[\{x\}]\mu(dx)$.}

\noindent Kita selalu mempunyai
$\hat\nu W[\{x\}]=\nu^*W[\{x\}]$ menurut definisi $\hat\nu$ (212C);
dan jika $\nu$ lengkap, maka $\hat\nu=\nu$, sehingga
$\lambda W=\int\nu W[\{x\}]\mu(dx)$.
}%end of proof of 252D

\leader{252E}{Akibat} Misalkan $(X,\Sigma,\mu)$ dan $(Y,\Tau,\nu)$
adalah ruang ukur, dengan produk c.l.d.\ $(X\times Y,\Lambda,\lambda)$.
Andaikan $\nu$ $\sigma$-hingga dan $\mu$
mempunyai himpunan terabaikan yang ditentukan secara lokal\cmmnt{ (213I)}.
Maka, jika $f$ adalah fungsi bernilai real yang terukur terhadap $\Lambda$
dan didefinisikan pada suatu subhimpunan $X\times Y$, fungsi
$y\mapsto f(x,y)$ terukur secara virtual terhadap $\nu$ untuk
$\mu$-hampir setiap $x\in X$.

\proof{ Misalkan $\tilde f$ adalah perluasan yang terukur terhadap
$\Lambda$ dari $f$ menjadi fungsi bernilai real yang didefinisikan di seluruh
$X\times Y$ (121I), dan tetapkan
$\tilde f_x(y)=\tilde f(x,y)$ untuk semua $x\in X$, $y\in Y$,

\Centerline{$D=\{x:x\in X,\,\tilde f_x$ terukur secara virtual terhadap
$\nu$ itu$\}$.}

Jika $G\in\Sigma$ dan $\mu G<\infty$, maka $G\setminus D$ terabaikan.
\Prf\ Misalkan $\sequencen{Y_n}$ adalah barisan tak-menurun himpunan
berukuran berhingga yang menutupi $Y$, dan tetapkan

$$\eqalign{\tilde f_n(x,y)
&=\tilde f(x,y)\text{ jika }x\in G,\,y\in Y_n
  \text{ dan }|\tilde f(x,y)|\le n,\cr
&=0\text{ di titik-titik lain }(x,y)\in X\times Y.\cr}$$

\noindent Maka setiap $\tilde f_n$ terintegralkan terhadap $\lambda$,
karena terbatas, terukur terhadap $\Lambda$, dan nol di luar
$G\times Y_n$. Akibatnya, dengan menetapkan
$\tilde f_{nx}(y)=\tilde f_n(x,y)$,

\Centerline{$\int(\int\tilde f_{nx}d\nu)\mu(dx)$ terdefinisi dan sama dengan
$\int\tilde f_nd\lambda$.}

\noindent Namun ini tentu berarti bahwa $\tilde f_{nx}$ terintegralkan
terhadap $\nu$, sehingga terukur secara virtual terhadap $\nu$, untuk
hampir setiap $x\in X$. Tetapkan

\Centerline{$D_n=\{x:x\in X,\,\tilde f_{nx}$ terukur secara virtual
terhadap $\nu$ itu$\}$;}

\noindent maka setiap $D_n$ koterabaikan terhadap $\mu$, sehingga
$\bigcap_{n\in\Bbb N}D_n$ koterabaikan. Tetapi untuk setiap
$x\in G\cap\bigcap_{n\in\Bbb N}D_n$,
$\tilde f_x=\lim_{n\to\infty}\tilde f_{nx}$ terukur secara virtual
terhadap $\nu$. Jadi
$G\setminus D\subseteq X\setminus\bigcap_{n\in\Bbb N}D_n$ terabaikan.\ \Qed

Hal ini berlaku setiap kali $\mu G<\infty$. Karena $G$ sebarang dan
$\mu$ mempunyai himpunan terabaikan yang ditentukan secara lokal, $D$
koterabaikan. Namun, untuk setiap $x\in D$, $y\mapsto f(x,y)$ merupakan
pembatasan dari $\tilde f_x$ dan harus terukur secara virtual terhadap
$\nu$.
}%end of proof of 252E

\leader{252F}{}\cmmnt{ Sebagai akibat lebih lanjut, kita dapat memperoleh
beberapa informasi berguna tentang ukuran produk c.l.d.\ untuk ruang ukur
sebarang.

\medskip

\noindent}{\bf Akibat} Misalkan $(X,\Sigma,\mu)$ dan $(Y,\Tau,\nu)$
adalah dua ruang ukur, $\lambda$ ukuran produk c.l.d.\ pada $X\times Y$,
dan $\Lambda$ domainnya. Misalkan $W\in\Lambda$ sedemikian sehingga
penampang vertikal $W[\{x\}]$ terabaikan terhadap $\nu$ untuk
$\mu$-hampir setiap $x\in X$. Maka $\lambda W=0$.

\proof{ Ambil $E\in\Sigma$, $F\in\Tau$ yang berukuran berhingga.
Misalkan $\lambda_{E\times F}$ adalah ukuran subruang pada $E\times F$.
Sekali lagi menurut 251Q(ii-$\alpha$), ukuran ini tepat merupakan produk
dari ukuran-ukuran subruang $\mu_E$ dan $\nu_F$. Kita tahu bahwa
$W\cap(E\times F)$ terukur oleh $\lambda_{E\times F}$. Pada saat yang
sama, penampang vertikal
$(W\cap(E\times F))[\{x\}]=W[\{x\}]\cap F$ terabaikan terhadap $\nu_F$
untuk $\mu_E$-hampir setiap $x\in E$. Dengan menerapkan 252B pada
$\mu_E$, $\nu_F$, dan $\chi(W\cap(E\times F))$,

\Centerline{$\lambda(W\cap(E\times F))
=\lambda_{E\times F}(W\cap(E\times F))
=\int_E\nu_F(W[\{x\}]\cap F)\mu_E(dx)=0$.}

\noindent Namun, dengan melihat definisi dalam 251F, kita melihat bahwa
ini berarti $\lambda W=0$, sebagaimana dinyatakan.
}%end of proof of 252F
\leader{252G}{}\cmmnt{ Teorema 252B dan akibat-akibatnya bergantung pada
ukuran faktor $\mu$ dan $\nu$ yang termasuk dalam kelas-kelas terbatas.
Ada suatu hasil parsial yang berlaku bagi semua ukuran produk c.l.d., sebagai
berikut.

\medskip

\noindent}{\bf Teorema Tonelli} Misalkan $(X,\Sigma,\mu)$ dan
$(Y,\Tau,\nu)$ ruang ukur, dan $(X\times Y,\Lambda,\lambda)$ produk c.l.d.
keduanya. Misalkan $f$ suatu fungsi yang terukur terhadap $\Lambda$, bernilai
$[-\infty,\infty]$, dan didefinisikan pada suatu anggota $\Lambda$, serta
andaikan bahwa salah satu dari
$\iint|f(x,y)|\mu(dx)\nu(dy)$ atau $\iint|f(x,y)|\nu(dy)\mu(dx)$ ada dalam
$\Bbb R$. Maka $f$ terintegralkan terhadap $\lambda$.

\proof{ Karena konstruksi ukuran produk bersifat simetris terhadap kedua
faktor, cukup kita tinjau kasus ketika $\iint
|f(x,y)|\nu(dy)\mu(dx)$ terdefinisi dan berhingga, sebab gagasan yang sama
tentu juga menangani kasus lainnya.

\medskip

{\bf (a)} Langkah pertama adalah memeriksa bahwa $f$ didefinisikan dan
berhingga hampir di mana-mana terhadap $\lambda$. \Prf\ Tetapkan
$W=\{(x,y):(x,y)\in\dom f,\,f(x,y)$ berhingga$\}$. Maka $W\in\Lambda$.
Hipotesis

\Centerline{$\iint|f(x,y)|\nu(dy)\mu(dx)$ terdefinisi dan berhingga}

\noindent mencakup pernyataan bahwa

\Centerline{$\int|f(x,y)|\nu(dy)$ terdefinisi dan berhingga
untuk $\mu$-hampir setiap $x$,}

\noindent yang menyiratkan bahwa

\Centerline{untuk $\mu$-hampir setiap $x$, $f(x,y)$ didefinisikan dan
berhingga untuk $\nu$-hampir setiap $y$;}

\noindent yaitu, bahwa

\Centerline{untuk $\mu$-hampir setiap $x$, $W[\{x\}]$
koterabaikan terhadap $\nu$.}

\noindent Tetapi menurut 252F hal ini menyiratkan bahwa
$(X\times Y)\setminus W$ terabaikan terhadap $\lambda$, sebagaimana
diperlukan.\ \Qed

\medskip

{\bf (b)} Misalkan $h$ sebarang fungsi sederhana-$\lambda$ yang nonnegatif
sedemikian sehingga $h\le|f|\,\,\lambda$-hampir di mana-mana.
Maka
\ifdim\pagewidth>467pt$\int h$ tidak mungkin lebih besar daripada
$\iint|f(x,y)|\nu(dy)\mu(dx)$.
\else $\int h$ paling besar $\iint|f(x,y)|\nu(dy)\mu(dx)$.\fi
\Prf\ Tetapkan

\Centerline{$W=\{(x,y):(x,y)\in\dom f,\,h(x,y)\le|f(x,y)|\}$,
\quad$h'=h\times\chi W$;}

\noindent maka $h'$ suatu fungsi sederhana dan $h'\eae h$. Nyatakan $h'$
sebagai $\sum_{i=0}^na_i\chi W_i$, dengan $a_i\ge 0$ dan
$\lambda W_i<\infty$ untuk setiap $i$.
Misalkan $\epsilon>0$. Untuk setiap $i\le n$ terdapat $E_i\in\Sigma$,
$F_i\in\Tau$ sedemikian sehingga $\mu E_i<\infty$, $\nu F_i<\infty$ dan
$\lambda(W_i\cap(E_i\times F_i))\ge\lambda W_i-\epsilon$. Tetapkan
$E=\bigcup_{i\le n}E_i$ dan $F=\bigcup_{i\le n}F_i$. Tinjau ukuran-ukuran
subruang $\mu_E$ dan $\nu_F$ serta produk keduanya
$\lambda_{E\times F}$ pada $E\times F$; maka $\lambda_{E\times F}$ adalah
ukuran subruang pada $E\times F$ yang didefinisikan dari $\lambda$
(sekali lagi 251Q(ii-$\alpha$)).
Karena itu, dengan menerapkan 252B pada produk $\mu_E\times\nu_F$,

\Centerline{$\int_{E\times F}h'\,d\lambda
=\int_{E\times F}h'\,d\lambda_{E\times F}
=\int_E\int_Fh'(x,y)\nu_F(dy)\mu_E(dx)$.}

Untuk sebarang $x$, kita mengetahui bahwa $h'(x,y)\le|f(x,y)|$ kapan pun
$f(x,y)$ didefinisikan. Jadi kita dapat memastikan bahwa

\Centerline{$\int_Fh'(x,y)\nu_F(dy)
=\int h'(x,y)\chi F(y)\nu(dy)
\le\int|f(x,y)|\nu(dy)$}

\noindent setidaknya kapan pun $\int_Fh'(x,y)\nu_F(dy)$ dan
$\int|f(x,y)|\nu(dy)$ keduanya terdefinisi, dan ini berlaku untuk hampir
setiap $x\in E$. Akibatnya

$$\eqalign{\int_{E\times F}h'\,d\lambda
&=\int_E\int_Fh'(x,y)\nu_F(dy)\mu_E(dx)\cr
&\le\int_E\int|f(x,y)|\nu(dy)\mu(dx)
\le\iint|f(x,y)|\nu(dy)\mu(dx).\cr}$$

Di sisi lain,

$$\eqalign{\int h'\,d\lambda-\int_{E\times F}h'\,d\lambda
&=\sum_{i=0}^na_i\lambda(W_i\setminus(E\times F))\cr
&\le\sum_{i=0}^na_i\lambda(W_i\setminus(E_i\times F_i))
\le\epsilon\sum_{i=0}^na_i.\cr}$$

\noindent Jadi

\Centerline{$\int h\,d\lambda
=\int h'd\lambda
\le\iint|f(x,y)|\nu(dy)\mu(dx)+\epsilon\sum_{i=0}^na_i$.}

\noindent Karena $\epsilon$ sebarang,
$\int h\,d\lambda
\le\iint|f(x,y)|\nu(dy)\mu(dx)$, sebagaimana dinyatakan.\ \Qed

\medskip

{\bf (c)} Hal ini berlaku kapan pun $h$ merupakan fungsi sederhana
terhadap $\lambda$ yang kurang dari atau sama dengan
$|f|\,\,\lambda$-hampir di mana-mana. Tetapi $|f|$ terukur terhadap
$\Lambda$ dan $\lambda$ semihingga (251Ic), sehingga ini cukup untuk
memastikan bahwa $|f|$ terintegralkan terhadap $\lambda$ (213B), yang
(karena $f$ diandaikan terukur terhadap $\Lambda$) pada gilirannya
menyiratkan bahwa $f$ terintegralkan terhadap $\lambda$.
}%end of proof of 252G

\leader{252H}{\bf Akibat} Misalkan $(X,\Sigma,\mu)$ dan $(Y,\Tau,\nu)$
ruang ukur $\sigma$-hingga, $\lambda$ ukuran produk c.l.d. pada
$X\times Y$, dan $\Lambda$ domainnya.

(a) Misalkan $f$ suatu fungsi yang terukur terhadap $\Lambda$, bernilai
$[-\infty,\infty]$, dan didefinisikan pada suatu anggota $\Lambda$. Maka
jika salah satu dari

\Centerline{$\int_{X\times Y}|f(x,y)|\lambda(d(x,y))$,
\quad$\int_Y\int_X|f(x,y)|\mu(dx)\nu(dy)$,
\quad$\int_X\int_Y|f(x,y)|\nu(dy)\mu(dx)$}

\noindent ada dalam $\Bbb R$, demikian pula dua lainnya, dan dalam hal ini

\Centerline{$\int_{X\times Y}f(x,y)\lambda(d(x,y))
=\int_Y\int_Xf(x,y)\mu(dx)\nu(dy)
=\int_X\int_Yf(x,y)\nu(dy)\mu(dx)$.}

(b)\dvAnew{2013} Misalkan $f$ suatu fungsi yang terukur terhadap $\Lambda$,
bernilai $[0,\infty]$, dan didefinisikan pada suatu anggota $\Lambda$.
Maka

\Centerline{$\int_{X\times Y}f(x,y)\lambda(d(x,y))
=\int_Y\int_Xf(x,y)\mu(dx)\nu(dy)
=\int_X\int_Yf(x,y)\nu(dy)\mu(dx)$}

\noindent dalam arti bahwa jika salah satu integral terdefinisi dalam
$[0,\infty]$, demikian pula dua lainnya, dan ketiganya kemudian sama.

\proof{{\bf (a)(i)} Andaikan bahwa $\int|f|d\lambda$ berhingga.
Karena $\mu$ dan $\nu$ keduanya $\sigma$-hingga, 252B memberi tahu kita
bahwa

\Centerline{$\iint|f(x,y)|\mu(dx)\nu(dy)$,
\quad$\iint|f(x,y)|\nu(dy)\mu(dx)$}

\noindent keduanya ada dan sama dengan $\int|f|d\lambda$, sedangkan

\Centerline{$\iint f(x,y)\mu(dx)\nu(dy)$,
\quad$\iint f(x,y)\nu(dy)\mu(dx)$}

\noindent keduanya ada dan sama dengan $\int fd\lambda$.

\medskip

\quad{\bf (ii)} Sekarang andaikan bahwa
$\iint|f(x,y)|\nu(dy)\mu(dx)$ ada dalam $\Bbb R$. Maka 252G memberi tahu
kita bahwa $|f|$ terintegralkan terhadap $\lambda$, sehingga kita dapat
menggunakan (i) untuk menyelesaikan argumen. Dengan menukar koordinat,
argumen yang sama berlaku jika $\iint|f(x,y)|\mu(dx)\nu(dy)$ ada dalam
$\Bbb R$.

\medskip

{\bf (b)(i)} Jika $\int fd\lambda$ terdefinisi, hasilnya langsung
mengikuti dari 252B.

\medskip

\quad{\bf (ii)} Andaikan bahwa $\iint_Xf(x,y)\nu(dy)\mu(dx)$ terdefinisi.
Seperti dalam bagian (a) bukti 252G, tetapi kali ini dengan menetapkan
$W=\dom f$, kita melihat bahwa $W\in\Lambda$ dan $W[\{x\}]$
koterabaikan terhadap $\nu$ untuk $\mu$-hampir setiap $x$,
sehingga $W$ koterabaikan terhadap $\lambda$. Karena $f$ nonnegatif,
terukur terhadap $\Lambda$, dan didefinisikan hampir di mana-mana,
$\int fd\lambda$ terdefinisi dalam $[0,\infty]$ dan kita berada dalam
kasus (i).
}%end of proof of 252H

\leader{252I}{Akibat} Misalkan $(X,\Sigma,\mu)$ dan $(Y,\Tau,\nu)$
ruang ukur, $\lambda$ ukuran produk c.l.d. pada $X\times
Y$, dan $\Lambda$ domainnya. Ambil $W\in\Lambda$. Jika salah satu dari
integral

\Centerline{$\int\mu^*W^{-1}[\{y\}]\nu(dy)$,
\quad$\int\nu^*W[\{x\}]\mu(dx)$}

\noindent ada dan berhingga, maka $\lambda W<\infty$.

\proof{ Terapkan 252G dengan $f=\chi W$, dengan mengingat bahwa

\Centerline{$\mu^*W^{-1}[\{y\}]=\int\chi W(x,y)\mu(dx)$,
\quad$\nu^*W[\{x\}]=\int\chi W(x,y)\nu(dy)$}

\noindent kapan pun integral-integral tersebut terdefinisi, seperti dalam
bukti 252D.
}%end of proof of 252I

\cmmnt{
\leader{252J}{Catatan} 252H merupakan bentuk dasar teorema Fubini; bukan
suatu kebetulan bahwa kebanyakan penulis menghindari ruang yang tidak
$\sigma$-hingga dalam konteks ini. Dua contoh berikut memperlihatkan
beberapa kesulitan yang dapat timbul jika kita meninggalkan wilayah yang
lazim berupa fungsi-fungsi yang kurang lebih terukur Borel pada ruang
$\sigma$-hingga. Contoh pertama bersifat klasik.
}%end of comment

\leader{252K}{Contoh} Misalkan $(X,\Sigma,\mu)$ adalah $[0,1]$ dengan
ukuran Lebesgue, dan $(Y,\Tau,\nu)$ adalah $[0,1]$ dengan ukuran cacah.
\cmmnt{

\spheader 252Ka} Tinjau himpunan

\Centerline{$W=\{(t,t):t\in[0,1]\}\subseteq X\times Y$.}

\cmmnt{
\noindent Kita mengamati bahwa $W$ dapat dinyatakan sebagai

\Centerline{$\bigcap_{n\in\Bbb N}\bigcup_{k=0}^n
[\Bover{k}{n+1},\Bover{k+1}{n+1}]\times[\Bover{k}{n+1},
\Bover{k+1}{n+1}]\in\Sigma\tensorhat\Tau$.}

\noindent Jika kita meninjau penampang-penampang

\Centerline{$W^{-1}[\{t\}]=W[\{t\}]=\{t\}$}

\noindent untuk $t\in [0,1]$, kita memperoleh
}

\ifwithproofs
\Centerline{$\iint\chi W(x,y)\mu(dx)\nu(dy)
=\int\mu W^{-1}[\{y\}]\nu(dy)=\int 0\,\nu(dy)=0$,}

\Centerline{$\iint\chi W(x,y)\nu(dy)\mu(dx)
=\int\nu W[\{x\}]\mu(dx)=\int 1\,\mu(dx)=1$,}
\else
\Centerline{$\iint\chi W(x,y)\mu(dx)\nu(dy)
=0$,}

\Centerline{$\iint\chi W(x,y)\nu(dy)\mu(dx)
=1$,}
\fi

\noindent sehingga kedua integral berulang tersebut berbeda.
\cmmnt{Oleh karena itu, secara umum urutan integrasi berulang tidak dapat
dibalik, bahkan untuk suatu fungsi terukur nonnegatif ketika kedua integral
berulang ada dan berhingga.
}%end of comment

\cmmnt{
\spheader 252Kb Karena himpunan $W$ pada bagian (a) sebenarnya termasuk
dalam $\Sigma\tensorhat\Tau$, kita mengetahui bahwa himpunan itu diukur oleh
ukuran produk c.l.d.\ $\lambda$, dan 252F (yang diterapkan dengan koordinat
dibalik) memberi tahu kita bahwa $\lambda W=0$.
}%end of comment

\cmmnt{
\spheader 252Kc Sebenarnya mudah untuk memberikan uraian lengkap tentang
$\lambda$.

\medskip

\quad{\bf (i)} Pokoknya adalah bahwa suatu himpunan
$W\subseteq[0,1]\times[0,1]$ termasuk dalam domain $\Lambda$ dari
$\lambda$ jika dan hanya jika setiap penampang mendatar
$W^{-1}[\{y\}]$ terukur Lebesgue. \prooflet{\Prf\
($\alpha$) Jika $W\in\Lambda$, maka, untuk setiap $b\in[0,1]$,
$\lambda([0,1]\times\{b\})$ berhingga, sehingga
$W\cap([0,1]\times\{b\})$ merupakan himpunan berukuran hingga, dan

\Centerline{$\lambda(W\cap([0,1]\times\{b\}))
=\int\mu(W\cap([0,1]\times\{b\}))^{-1}[\{y\}]\nu(dy)
=\mu W^{-1}[\{b\}]$}

\noindent menurut 252D, karena $\mu$ bersifat $\sigma$-hingga, $\nu$
bersifat dapat dilokalkan secara ketat, lengkap, dan ditentukan secara
lokal, serta

$$\eqalign{(W\cap([0,1]\times\{b\}))^{-1}[\{y\}]
&=W^{-1}[\{b\}]\text{ jika }y=b,\cr
&=\emptyset\text{ selain itu.}\cr}$$

\noindent Karena $b$ sebarang, setiap penampang mendatar dari $W$ terukur.
($\beta$) Jika setiap penampang mendatar dari $W$ terukur, misalkan
$F\subseteq[0,1]$ sebarang himpunan berukuran hingga terhadap $\nu$; maka
$F$ berhingga, sehingga

\Centerline{$W\cap([0,1]\times F)
=\bigcup_{y\in F}W^{-1}[\{y\}]\times\{y\}
\in\Sigma\tensorhat\Tau\subseteq\Lambda$.}

\noindent Tetapi dari sini mengikuti bahwa $W$ sendiri termasuk dalam
$\Lambda$, menurut 251H.\ \Qed
}%end of prooflet

\medskip

\quad{\bf (ii)} Sekarang beberapa perhitungan yang sama menunjukkan bahwa
untuk setiap $W\in\Lambda$,

\Centerline{$\lambda W=\sum_{y\in [0,1]}\mu W^{-1}[\{y\}]$.}

\prooflet{\noindent\Prf\ Untuk sebarang himpunan berhingga
$F\subseteq[0,1]$,

$$\eqalign{\lambda(W\cap([0,1]\times F))
&=\int\mu(W\cap([0,1]\times F))^{-1}[\{y\}]\nu(dy)\cr
&=\int_F\mu W^{-1}[\{y\}]\nu(dy)
=\sum_{y\in F}\mu W^{-1}[\{y\}].\cr}$$

\noindent Jadi

\Centerline{$\lambda W=\sup_{F\subseteq[0,1]\text{ berhingga}}
\sum_{y\in F}\mu W^{-1}[\{y\}]=\sum_{y\in [0,1]}\mu W^{-1}[\{y\}]$.
\Qed}
}%end of prooflet
}%end of comment

\cmmnt{
\leader{252L}{Contoh} Untuk contoh kedua, saya beralih kepada suatu masalah
yang dapat timbul jika kita lalai memeriksa bahwa suatu fungsi terukur
sebagai fungsi dari dua variabel.

Misalkan $(X,\Sigma,\mu)=(Y,\Tau,\nu)$ adalah $\omega_1$, ordinal tak
terhitung pertama (2A1Fc), dengan ukuran terhitung-koterhitung (211R).
Tetapkan

\Centerline{$W=\{(\xi,\eta):\xi\le\eta<\omega_1\}\subseteq X\times Y$.}

\noindent Maka semua penampang mendatar
$W^{-1}[\{\eta\}]=\{\xi:\xi\le\eta\}$ terhitung, sehingga

\Centerline{$\int\mu W^{-1}[\{\eta\}]\nu(d\eta)
=\int 0\,\nu(d\eta)
=0$,}

\noindent sedangkan semua penampang tegak
$W[\{\xi\}]=\{\eta:\xi\le\eta<\omega_1\}$ koterhitung, sehingga

\Centerline{$\int\nu W[\{\xi\}]\mu(d\xi)
=\int 1\,\mu(d\xi)
=1$.}

\noindent Karena kedua integral berulang tersebut berbeda, keduanya tidak
mungkin sama dengan ukuran $W$, dan satu-satunya penyelesaian
adalah mengatakan bahwa $W$ tidak diukur oleh ukuran produk.
}%end of comment

\cmmnt{
\leader{252M}{Catatan} Jenis kesulitan ketiga dalam rumus

\Centerline{$\iint f(x,y)dxdy=\iint f(x,y)dydx$}

\noindent dapat timbul bahkan pada ruang probabilitas dengan fungsi
bernilai real yang terukur terhadap $\Sigma\tensorhat\Tau$ dan didefinisikan
di mana-mana jika kita lalai memeriksa bahwa $f$ terintegralkan terhadap
ukuran produk. Dalam 252H, kita memang memerlukan hipotesis bahwa salah satu
dari

\Centerline{$\int_{X\times Y}|f(x,y)|\lambda(d(x,y))$,
\quad$\int_Y\int_X|f(x,y)|\mu(dx)\nu(dy)$,
\quad$\int_X\int_Y|f(x,y)|\nu(dy)\mu(dx)$}

\noindent berhingga. Untuk contoh-contoh yang menunjukkan hal ini, lihat
252Xf dan 252Xg.
}%end of comment

\leader{252N}{Integrasi melalui himpunan ordinat I: Proposisi} Misalkan
$(X,\Sigma,\mu)$ suatu ruang ukur lengkap yang ditentukan secara lokal, dan
$\lambda$ ukuran produk c.l.d. pada $X\times\Bbb R$, dengan $\Bbb R$
diberi ukuran Lebesgue; tuliskan $\Lambda$ untuk domain $\lambda$.
Untuk sebarang fungsi bernilai $[0,\infty]$ $f$ yang didefinisikan pada
suatu subhimpunan koterabaikan dari $X$, tuliskan $\Omega_f$, $\Omega'_f$
untuk {\bf himpunan ordinat}

\Centerline{$\Omega_f=\{(x,a):x\in\dom f,\,0\le a\le f(x)\}
\subseteq X\times\Bbb R$,}

\Centerline{$\Omega'_f=\{(x,a):x\in\dom f,\,0\le a<f(x)\}
\subseteq X\times\Bbb R$.}

\noindent Maka

\Centerline{$\lambda\Omega_f=\lambda\Omega'_f=\int fd\mu$}

\noindent dalam arti bahwa jika salah satu di antaranya terdefinisi dalam
$[0,\infty]$, demikian pula dua lainnya, dan ketiganya sama.

\proof{{\bf (a)} Jika $\Omega_f\in\Lambda$, maka

\Centerline{$\int f(x)\mu(dx)
=\int\nu\{y:(x,y)\in\Omega_f\}\mu(dx)
=\lambda\Omega_f$}

\noindent menurut 252D, dengan menuliskan $\nu$ untuk ukuran Lebesgue,
karena $f$ didefinisikan hampir di mana-mana. Demikian pula, jika
$\Omega'_f\in\Lambda$,

\Centerline{$\int f(x)\mu(dx)
=\int\nu\{y:(x,y)\in\Omega'_f\}\mu(dx)
=\lambda\Omega'_f$.}

\medskip

{\bf (b)} Jika $\int fd\mu$ terdefinisi, maka $f$ terukur secara virtual
terhadap $\mu$, dan karenanya terukur (karena $\mu$ lengkap);
sekali lagi karena $\mu$ lengkap, $\dom f\in\Sigma$. Jadi

\Centerline{$\Omega'_f
=\bigcup_{q\in\Bbb Q,q>0}\{x:x\in\dom f,\,f(x)>q\}\times[0,q]$,}

\Centerline{$\Omega_f
=\bigcap_{n\ge 1}\bigcup_{q\in\Bbb Q,q>0}
  \{x:x\in\dom f,\,f(x)\ge q-\Bover1n\}\times[0,q]$}

\noindent termasuk dalam $\Lambda$, sehingga $\lambda\Omega_f$ dan
$\lambda\Omega'_f$ terdefinisi. Sekarang keduanya sama dengan
$\int fd\mu$, menurut (a).
}%end of proof of 252N

\leader{252O}{Integrasi melalui himpunan ordinat II: Proposisi} Misalkan
$(X,\Sigma,\mu)$ suatu ruang ukur, dan $f$ suatu fungsi nonnegatif yang
terukur secara virtual terhadap $\mu$ dan didefinisikan pada suatu
subhimpunan koterabaikan dari $X$. Maka

\Centerline{$\int fd\mu
=\int_0^{\infty}\mu^*\{x:x\in\dom f,\,f(x)\ge t\}dt
=\int_0^{\infty}\mu^*\{x:x\in\dom f,\,f(x)>t\}dt$}

\noindent dalam $[0,\infty]$, dengan integral-integral $\int\ldots dt$
diambil terhadap ukuran Lebesgue.

\proof{ Melengkapi $\mu$ tidak mengubah integral $f$ ataupun ukuran luar
$\mu^*$ (212Fb, 212Ea), sehingga kita boleh mengandaikan bahwa $\mu$
lengkap; dalam hal ini $\dom f$ dan $f$ akan terukur.
Untuk $n$, $k\in\Bbb N$ tetapkan
$E_{nk}=\{x:x\in\dom f,\,f(x)>2^{-n}k\}$,
$g_n(x)=2^{-n}\sum_{k=1}^{4^n}\chi E_{nk}$. Maka $\sequencen{g_n}$ adalah
barisan tak menurun dari fungsi-fungsi terukur yang konvergen ke $f$ pada
setiap titik $\dom f$, sehingga
$\int fd\mu=\lim_{n\to\infty}\int g_nd\mu$ dan
$\mu\{x:f(x)>t\}=\lim_{n\to\infty}\mu\{x:g_n(x)>t\}$ untuk setiap
$t\ge 0$; akibatnya

\Centerline{$\int_0^{\infty}\mu\{x:f(x)>t\}dt
=\lim_{n\to\infty}\int_0^{\infty}\mu\{x:g_n(x)>t\}dt$.}

\noindent Di sisi lain,
$\mu\{x:g_n(x)>t\}=\mu E_{nk}$ jika $1\le k\le 4^n$ dan
$2^{-n}(k-1)\le t<2^{-n}k$, dan $0$ jika $t\ge 2^n$, sehingga

\Centerline{$\int_0^{\infty}\mu\{x:g_n(x)>t\}dt$
$=\sum_{k=1}^{4^n}2^{-n}\mu E_{nk}$
$=\int g_n\,d\mu$,}

\noindent untuk setiap $n\in\Bbb N$. Jadi
$\int_0^{\infty}\mu\{x:f(x)>t\}dt=\int fd\mu$.

Sekarang $\mu\{x:f(x)\ge t\}=\mu\{x:f(x)>t\}$ untuk hampir setiap $t$.
\Prf\ Tetapkan $C=\{t:\mu\{x:f(x)>t\}<\infty\}$,
$h(t)=\mu\{x:f(x)>t\}$ untuk $t\in C$. Jika $C$ tidak kosong,
$h:C\to\coint{0,\infty}$ monoton, sehingga kontinu hampir di mana-mana
dalam $C$ (222A). Tetapi pada setiap titik
$C\setminus\{\inf C\}$ tempat $h$ kontinu,

\Centerline{$\mu\{x:f(x)\ge t\}=\lim_{s\uparrow t}\mu\{x:f(x)>s\}
=\mu\{x:f(x)>t\}$.}

\noindent Jadi kita memperoleh hasilnya, karena
$\mu\{x:f(x)\ge t\}=\mu\{x:f(x)>t\}=\infty$ untuk sebarang
$t\in\coint{0,\infty}\setminus C$.\ \Qed

Dengan demikian $\int_0^{\infty}\mu\{x:f(x)\ge t\}dt$ juga sama dengan
$\int fd\mu$.
}%end of proof of 252O

\leader{*252P}{}\cmmnt{ Jika kita menerapkan gagasan-gagasan 252B secara
lengkap pada fungsi-fungsi yang terukur terhadap
$\Sigma\tensorhat\Tau$, kita memperoleh hasil berikut, yang kadang-kadang
berguna.

\medskip

\noindent}{\bf Proposisi} Misalkan $(X,\Sigma,\mu)$ suatu ruang ukur, dan
$(Y,\Tau,\nu)$ suatu ruang ukur $\sigma$-hingga. Maka untuk sebarang fungsi
yang terukur terhadap $\Sigma\tensorhat\Tau$ dengan bentuk
$f:X\times Y\to[0,\infty]$,
$x\mapsto\int f(x,y)\nu(dy):X\to[0,\infty]$ terukur terhadap $\Sigma$.
Jika $\mu$ semihingga,
$\iint f(x,y)\nu(dy)\mu(dx)\discretionary{}{}{}=\int fd\lambda$,
dengan $\lambda$ ukuran produk c.l.d. pada $X\times Y$.

\proof{{\bf (a)} Misalkan $\sequencen{Y_n}$ suatu barisan tak menurun dari
subhimpunan-subhimpunan $Y$ yang berukuran hingga dan gabungannya adalah
$Y$. Tetapkan

$$\eqalign{\Cal A
&=\{W:W\subseteq X\times Y,\,W[\{x\}]\in\Tau\text{ untuk setiap }x\in X,\cr
&\qquad\qquad\qquad\qquad
x\mapsto\nu(Y_n\cap W[\{x\}])\text{ terukur terhadap }\Sigma\text{ untuk
setiap }n\in\Bbb N\}.\cr}$$

\noindent Maka $\Cal A$ merupakan suatu kelas Dynkin dari subhimpunan
$X\times Y$ yang mencakup $\{E\times F:E\in\Sigma,\,F\in\Tau\}$, sehingga
mencakup $\Sigma\tensorhat\Tau$, sekali lagi menurut Teorema Kelas Monoton
(136B).

Ini berarti bahwa jika $W\in\Sigma\tensorhat\Tau$, maka

\Centerline{$\nu W[\{x\}]=\sup_{n\in\Bbb N}\nu(Y_n\cap W[\{x\}])$}

\noindent terdefinisi untuk setiap $x\in X$ dan terukur terhadap $\Sigma$
sebagai fungsi dari $x$.

\medskip

{\bf (b)} Sekarang, untuk $n$, $k\in\Bbb N$, tetapkan

\Centerline{$W_{nk}=\{(x,y):f(x,y)\ge 2^{-n}k\}$,
\quad$g_n=\sum_{k=1}^{4^n}2^{-n}\chi W_{nk}$.}

\noindent Lalu jika kita menetapkan

\Centerline{$h_n(x)
=\int g_n(x,y)\nu(dy)=\sum_{k=1}^{4^n}2^{-n}\nu W_{nk}[\{x\}]$}

\noindent untuk $n\in\Bbb N$ dan $x\in X$,
$h_n:X\to[0,\infty]$ terukur terhadap $\Sigma$, dan

\Centerline{$\lim_{n\to\infty}h_n(x)
=\int(\lim_{n\to\infty}g_n(x,y))\nu(dy)=\int f(x,y)\nu(dy)$}

\noindent untuk setiap $x$, karena $\sequencen{g_n(x,y)}$ adalah suatu
barisan tak menurun dengan limit $f(x,y)$ untuk semua $x\in X$, $y\in Y$.
Jadi $x\mapsto\int f(x,y)\nu(dy)$ didefinisikan di seluruh $X$ dan terukur
terhadap $\Sigma$.

\medskip

{\bf (c)} Jika $E\subseteq X$ terukur dan berukuran hingga, maka
$\int_E\int f(x,y)\nu(dy)\mu(dx)=\int_{E\times Y}fd\lambda$, dengan
menerapkan 252B pada produk ukuran subruang $\mu_E$ dan $\nu$ (dan
menggunakan 251Q untuk memeriksa bahwa produk $\mu_E$ dan $\nu$ adalah
ukuran subruang pada $E\times Y$). Sekarang jika $\lambda W$ terdefinisi
dan berhingga, harus ada suatu barisan tak menurun $\sequencen{E_n}$ dari
subhimpunan-subhimpunan $X$ yang berukuran hingga sedemikian sehingga
$\lambda W=\sup_{n\in\Bbb N}\lambda(W\cap(E_n\times Y))$, sehingga
$W\setminus\bigcup_{n\in\Bbb N}(E_n\times Y)$ terabaikan, dan

$$\eqalignno{\int_Wfd\lambda
&=\lim_{n\to\infty}\int_{W\cap(E_n\times Y)}fd\lambda\cr
\displaycause{menurut teorema B.Levi yang diterapkan pada
$\sequencen{f\times\chi(W\cap(E_n\times Y))}$}
&\le\lim_{n\to\infty}\int_{E_n\times Y}fd\lambda
=\lim_{n\to\infty}\int_{E_n}\int f(x,y)\nu(dy)\mu(dx)\cr
&\le\iint f(x,y)\nu(dy)\mu(dx).\cr}$$

\noindent Sekali lagi menurut 213B,

\Centerline{$\int fd\lambda
=\sup_{\lambda W<\infty}\int_Wfd\lambda
\le\iint f(x,y)\nu(dy)\mu(dx)$.}

\noindent Tetapi juga, jika $\mu$ semihingga,

\Centerline{$\iint f(x,y)\nu(dy)\mu(dx)
=\sup_{\mu E<\infty}\int_E\int f(x,y)\nu(dy)\mu(dx)
\le\int fd\lambda$,}

\noindent sehingga $\int fd\lambda=\iint f(x,y)\nu(dy)\mu(dx)$,
sebagaimana dinyatakan.
}%end of proof of 252P

\leader{252Q}{Volume bola}\dvro{ Ukuran Lebesgue bola
satuan di $\BbbR^r$ adalah}
{ Sekarang kita telah memiliki semua
perangkat pokok untuk melakukan sebuah perhitungan kecil yang, saya kira,
layak diketahui sekadar sebagai pengetahuan umum: volume bola satuan
$\{x:\|x\|\le 1\}=\{(\xi_1,\ldots,\xi_r):\sum_{i=1}^r\xi_i^2\le 1\}$ di
$\BbbR^r$. Sebenarnya, dari sudut pandang teoretis,
saya rasa kita hampir dapat menamainya saja $\beta_r$ dan berhenti
sampai di situ; tetapi karena terdapat rumus umum yang dinyatakan dengan
$\beta_2=\pi$ dan faktorial, rasanya memalukan jika tidak disajikan. Perhitungan
ini sama sekali tidak berkaitan dengan integrasi Lebesgue, dan saya dapat
menyingkirkannya sebagai sekadar kalkulus lanjut; tetapi karena kini hanya
sebagian kecil matematikawan diajari kalkulus sampai taraf ini dengan ketelitian
yang memadai sebelum diperkenalkan kepada integral Lebesgue, saya tidak ragu
bahwa banyak pembaca, seperti saya sendiri, melewatkan beberapa seluk-beluk
yang terlibat. Karena itu saya menyediakan ruang untuk menguraikan rinciannya
dengan gaya yang digunakan di bagian lain jilid ini, seraya menyadari bahwa
perangkat yang dipakai jauh lebih rumit daripada yang sungguh diperlukan untuk
hasil ini.

\header{252Qa}{\bf (a)} Fakta dasar pertama yang kita perlukan ialah bahwa,
untuk setiap $n\ge 1$,

$$\eqalign{I_n=\int_{-\pi/2}^{\pi/2}\cos^nt\,dt
&=\Bover{(2k)!}{(2^kk!)^2}\pi\text{ jika }n=2k\text{ genap},\cr
&=2\Bover{(2^kk!)^2}{(2k+1)!}\text{ jika }n=2k+1\text{ ganjil}.\cr}$$

\prooflet{\noindent \Prf\  Untuk $n=0$, tentu saja,

\Centerline{$I_0=\int_{-\pi/2}^{\pi/2}1\,dt=\pi
=\Bover{0!}{(2^00!)^2}\pi$,}

\noindent sedangkan untuk $n=1$ kita memperoleh

\Centerline{$I_1=\sin\bover{\pi}2-\sin(-\bover{\pi}2)=2
=2\Bover{(2^00!)^2}{1!}$,}

\noindent dengan menggunakan Teorema Dasar Kalkulus (225L) dan fakta
bahwa $\sin'=\cos$ terbatas. Untuk langkah induksi menuju $n+1\ge 2$, kita
dapat menggunakan integrasi parsial (225F):

$$\eqalign{I_{n+1}
&=\int_{-\pi/2}^{\pi/2}\cos t\cos^nt\,dt\cr
&=\sin\Bover{\pi}2\cos^n\Bover{\pi}2
   -\sin(-\Bover{\pi}2)\cos^n(-\Bover{\pi}2)
   +\int_{-\pi/2}^{\pi/2}\sin t\cdot n\cos^{n-1}t\cdot\sin t\,dt\cr
&=n\int_{-\pi/2}^{\pi/2}(1-\cos^2t)\cos^{n-1}t\,dt
=n(I_{n-1}-I_{n+1}),\cr}$$

\noindent sehingga $I_{n+1}=\Bover{n}{n+1}I_{n-1}$.
Sekarang rumus-rumus yang diberikan mengikuti melalui induksi sederhana.\ \Qed}

\spheader 252Qb Hasil berikutnya ialah bahwa, untuk setiap $n\in\Bbb N$
dan setiap $a\ge 0$,

\Centerline{$\int_{-a}^a(a^2-s^2)^{n/2}ds=I_{n+1}a^{n+1}$.}

\prooflet{\noindent\Prf\ Tentu saja ini merupakan integrasi dengan
substitusi; tetapi singularitas integran di $s=\pm a$ sedikit
memperumit persoalan. Saya kemukakan argumen berikut. Jika $a=0$, hasilnya
sepele; ambil $a>0$. Untuk $-a\le b\le a$, tetapkan
$F(b)=\int_{-a}^b(a^2-s^2)^{n/2}ds$. Karena integrannya kontinu,
$F'(b)$ ada dan sama dengan $(a^2-b^2)^{n/2}$ untuk
$-a<b<a$ (222H). Tetapkan $G(t)=F(a\sin t)$; maka $G$ kontinu dan

\Centerline{$G'(t)=aF'(a\sin t)\cos t=a^{n+1}\cos^{n+1}t$}

\noindent untuk
$-\bover{\pi}2<t<\bover{\pi}2$. Akibatnya

$$\eqalignno{\int_{-a}^a(a^2-s^2)^{n/2}ds
&=F(a)-F(-a)
=G(\Bover{\pi}2)-G(-\Bover{\pi}2)
=\int_{-\pi/2}^{\pi/2}G'(t)dt\cr
\noalign{\noindent (menurut 225L, seperti sebelumnya)}
&=a^{n+1}I_{n+1},\cr}$$

\noindent sebagaimana diperlukan.\ \Qed}

\spheader 252Qc Akhirnya sekarang kita siap membahas bola-bola
$B_r=\{x:x\in\BbbR^r,\,\|x\|\le 1\}$. Misalkan $\mu_r$ ukuran Lebesgue
pada $\BbbR^r$, dan tetapkan $\beta_r=I_1I_2\ldots I_r$ untuk $r\ge 1$.
Saya nyatakan bahwa, dengan menuliskan

\Centerline{$B_r(a)=\{x:x\in\BbbR^r,\,\|x\|\le a\}$,}

\noindent kita memperoleh $\mu_r(B_r(a))=\beta_ra^r$ untuk setiap $a\ge 0$.
\prooflet{\Prf\ Lakukan induksi pada $r$. Untuk $r=1$ kita mempunyai
$\beta_1=2$, $B_1(a)=[-a,a]$, sehingga hasilnya sepele. Untuk langkah
induksi menuju $r+1$, kita memperoleh

$$\eqalignno{\mu_{r+1}B_{r+1}(a)
&=\int\mu_r\{x:(x,t)\in B_{r+1}(a)\}dt \cr
\noalign{\noindent (dengan menggabungkan 251N dan 252D, serta menggunakan fakta
bahwa $B_{r+1}(a)$ tertutup, sehingga terukur)}
&=\int_{-a}^a\mu_rB_r(\sqrt{a^2-t^2})dt\cr
\noalign{\noindent (karena $(x,t)\in B_{r+1}(a)$ jika dan hanya jika $|t|\le a$ dan
$\|x\|\le\sqrt{a^2-t^2}$)}
&=\int_{-a}^a\beta_r(a^2-t^2)^{r/2}dt\cr
\noalign{\noindent (menurut hipotesis induksi)}
&=\beta_ra^{r+1}I_{r+1}\cr
\noalign{\noindent (menurut (b) di atas)}
&=\beta_{r+1}a^{r+1}\cr}$$

\noindent (menurut definisi $\beta_{r+1}$). Dengan demikian induksi
berlanjut.\ \Qed}

\spheader 252Qd Khususnya, ukuran Lebesgue berdimensi $r$ dari
bola berdimensi $r$, $B_r=B_r(1)$, hanyalah
$\beta_r=I_1\ldots I_r$. Sekarang induksi sederhana pada $k$ menunjukkan bahwa
}%end of \dvro

$$\eqalign{\beta_r
&=\Bover{1}{k!}\pi^k\text{ jika }r=2k\text{ genap},\cr
&=\Bover{2^{2k+1}k!}{(2k+1)!}\pi^k\text{ jika }r=2k+1\text{ ganjil}.\cr}$$

\cmmnt{
\spheader 252Qe Perhatikan bahwa pada bagian (c) dari pembuktian kita melihat bahwa
$\{x:x\in\BbbR^r,\,\|x\|\le a\}$ mempunyai ukuran $\beta_ra^r$ untuk setiap
$a\ge 0$.

Rumus-rumus di sini konsisten dengan penetapan $\beta_0=1$;
yang bersesuaian dengan pernyataan bahwa $\BbbR^0=\{\emptyset\}$, bahwa
$\mu_0\BbbR^0=1$ dan bahwa $B_0=\{\emptyset\}$. Menetapkan $\mu_0\BbbR^0$
sebagai $1$ juga konsisten dengan gagasan bahwa, mengikuti 251N, ukuran produk
$\mu_0\times\mu_r$ semestinya sesuai dengan $\mu_{0+r}$ pada $\BbbR^{0+r}$.
}%end of comment


\leader{252R}{Fungsi bernilai kompleks}\cmmnt{ Hasil-hasil 252B--252I
di atas mudah diterapkan pada fungsi bernilai kompleks, dengan
mempertimbangkan bagian riil dan bagian imajinernya. Secara khusus:

\medskip

}{\bf (a)} Misalkan $(X,\Sigma,\mu)$ dan $(Y,\Tau,\nu)$ adalah
ruang-ruang ukur, dan $\lambda$ ukuran produk c.l.d.\ pada $X\times
Y$. Andaikan bahwa $\nu$ $\sigma$-hingga dan bahwa $\mu$ dapat
dilokalkan secara ketat atau lengkap dan ditentukan secara lokal. Misalkan $f$
fungsi bernilai kompleks yang dapat diintegralkan terhadap $\lambda$. Maka
$\iint f(x,y)\nu(dy)\mu(dx)$ terdefinisi dan sama dengan $\int fd\lambda$.

\spheader 252Rb Misalkan $(X,\Sigma,\mu)$ dan $(Y,\Tau,\nu)$
adalah ruang-ruang ukur, $\lambda$ ukuran produk c.l.d.\ pada $X\times
Y$, dan $\Lambda$ domainnya.
Misalkan $f$ fungsi bernilai kompleks terukur terhadap $\Lambda$ yang didefinisikan
pada suatu anggota $\Lambda$, dan andaikan bahwa
$\iint|f(x,y)|\mu(dx)\nu(dy)$ atau $\iint|f(x,y)|\nu(dy)\mu(dx)$
terdefinisi dan berhingga. Maka $f$ dapat diintegralkan terhadap $\lambda$.

\spheader 252Rc Misalkan $(X,\Sigma,\mu)$ dan $(Y,\Tau,\nu)$ adalah
ruang-ruang ukur $\sigma$-hingga, $\lambda$ ukuran produk c.l.d.\
pada $X\times Y$, dan $\Lambda$ domainnya. Misalkan $f$ fungsi bernilai
kompleks terukur terhadap $\Lambda$ yang didefinisikan pada suatu anggota
$\Lambda$. Maka jika salah satu dari

\Centerline{$\int_{X\times Y}|f(x,y)|\lambda(d(x,y))$,
\quad$\int_Y\int_X|f(x,y)|\mu(dx)\nu(dy)$,
\quad$\int_X\int_Y|f(x,y)|\nu(dy)\mu(dx)$}

\noindent ada dalam $\Bbb R$, demikian pula dua yang lain, dan dalam hal ini

\Centerline{$\int_{X\times Y}f(x,y)\lambda(d(x,y))
=\int_Y\int_Xf(x,y)\mu(dx)\nu(dy)
=\int_X\int_Yf(x,y)\nu(dy)\mu(dx)$.}

\exercises{
\leader{252X}{Latihan dasar (a)}
%\spheader 252Xa
Misalkan $(X,\Sigma,\mu)$ dan $(Y,\Tau,\nu)$ adalah
ruang-ruang ukur, dan
$\lambda$ ukuran produk c.l.d.\ pada $X\times Y$. Misalkan $f$ fungsi
bernilai riil yang dapat diintegralkan terhadap $\lambda$ sedemikian sehingga
$\int_{E\times F}f=0$ bilamana $E\in\Sigma$, $F\in\Tau$, $\mu E<\infty$
dan $\nu F<\infty$. Tunjukkan bahwa $f=0\,\,\lambda$-hampir di mana-mana.
\Hint{gunakan 251Ie untuk menunjukkan bahwa $\int_Wf=0$ bilamana
$\lambda W<\infty$.}
%252B

\spheader 252Xb Misalkan $f$, $g:\Bbb R\to\Bbb R$ adalah dua fungsi
tak-menurun, dan $\mu_f$, $\mu_g$ ukuran-ukuran Lebesgue--Stieltjes yang
bersesuaian (lihat 114Xa). Tetapkan

\Centerline{$f(x^+)=\lim_{t\downarrow x}f(t)$,
\quad$f(x^-)=\lim_{t\uparrow x}f(t)$}

\noindent untuk setiap $x\in\Bbb R$, dan definisikan $g(x^+)$, $g(x^-)$
dengan cara serupa. Tunjukkan bahwa bilamana $a\le b$ dalam $\Bbb R$,

$$\eqalign{\int_{[a,b]}f(x^-)\mu_g(dx)&+\int_{[a,b]}g(x^+)\mu_f(dx)
=g(b^+)f(b^+)-g(a^-)f(a^-)\cr
&=\int_{[a,b]}\Bover12(f(x^-)+f(x^+))\mu_g(dx)
  +\int_{[a,b]}\Bover12(g(x^-)+g(x^+))\mu_f(dx).\cr}$$

\noindent \Hint{temukan dua ungkapan untuk
$(\mu_f\times\mu_g)\{(x,y):a\le x<y\le b\}$.}
%252C

\sqheader 252Xc Misalkan $(X,\Sigma,\mu)$ dan $(Y,\Tau,\nu)$ adalah
ruang-ruang ukur yang lengkap dan ditentukan secara lokal, $\lambda$ ukuran
produk c.l.d.\ pada $X\times Y$, dan $\Lambda$ domainnya. Andaikan bahwa
$A\subseteq X$ dan
$B\subseteq Y$. Tunjukkan bahwa $A\times B\in\Lambda$ jika dan hanya jika
$\mu A=0$ atau $\nu B=0$ atau $A\in\Sigma$ dan $B\in\Tau$.
\Hint{jika $B$ tidak terabaikan dan $A\times B\in\Lambda$, ambil $H$ sedemikian
sehingga $\nu H<\infty$ dan $B\cap H$ tidak terabaikan. Maka
$W=A\times(B\cap H)$ diukur oleh $\mu\times\nu_H$, dengan $\nu_H$ merupakan
ukuran subruang pada $H$. Sekarang terapkan 252D pada $\mu$, $\nu_H$ dan $W$
untuk melihat bahwa $A\in\Sigma$.}
%252D

\sqheader 252Xd Misalkan $(X_1,\Sigma_1,\mu_1)$,
$(X_2,\Sigma_2,\mu_2)$, $(X_3,\Sigma_3,\mu_3)$ adalah tiga ruang ukur
$\sigma$-hingga, dan $f$ fungsi bernilai riil yang didefinisikan hampir di
mana-mana pada $X_1\times X_2\times X_3$ dan terukur terhadap $\Lambda$,
dengan $\Lambda$ merupakan domain ukuran produk yang dideskripsikan dalam
251W atau 251Xg. Tunjukkan bahwa jika
$\iiint |f(x_1,x_2,x_3)|dx_1dx_2dx_3$ terdefinisi dalam $\Bbb R$, maka
$\iiint f(x_1,x_2,x_3)dx_2dx_3dx_1$ dan
$\iiint f(x_1,x_2,x_3)dx_3dx_1dx_2$
ada dan sama.
%252H

\spheader 252Xe Berikan contoh ruang-ruang ukur yang dapat dilokalkan secara
ketat $(X,\Sigma,\mu)$, $(Y,\Tau,\nu)$ dan suatu $W\in\Sigma\tensorhat\Tau$
sedemikian sehingga $x\mapsto\nu W[\{x\}]$ tidak terukur terhadap $\Sigma$.
\Hint{dalam 252Kb, coba ambil $Y$ sebagai subhimpunan sejati dari $[0,1]$.}
%252K

\sqheader 252Xf Tetapkan $f(x,y)=\sin(x-y)$ jika $0\le y\le x\le y+2\pi$, dan
$0$ di titik $(x,y)\in\BbbR^2$ lainnya. Tunjukkan bahwa
$\iint f(x,y)dx\,dy=0$ dan
$\iint f(x,y)dy\,dx=2\pi$, dengan semua integral diambil terhadap ukuran
Lebesgue.
%252M

\sqheader 252Xg Tetapkan $f(x,y)=\Bover{x^2-y^2}{(x^2+y^2)^2}$ untuk $x$,
$y\in\ocint{0,1}$. Tunjukkan bahwa
$\int_0^1\int_0^1f(x,y)dydx=\Bover{\pi}4$,
\discrversionA{\break}{}$\int_0^1\int_0^1f(x,y)dxdy=-\Bover{\pi}4$.
%252M

\sqheader 252Xh Misalkan $(X,\Sigma,\mu)$ dan $(Y,\Tau,\nu)$ adalah
ruang-ruang ukur, dan $f$ fungsi terukur terhadap
$\Sigma\tensorhat\Tau$ yang didefinisikan pada suatu subhimpunan
$X\times Y$. Tunjukkan bahwa $y\mapsto f(x,y)$ terukur terhadap $\Tau$
untuk setiap $x\in X$.
%252P

\spheader 252Xi Misalkan $r\ge 1$ bilangan bulat, dan tuliskan $\beta_r$ untuk
ukuran Lebesgue bola satuan di $\BbbR^r$. Tetapkan
$g_r(t)=r\beta_rt^{r-1}$ untuk $t\ge 0$, $\phi(x)=\|x\|$ untuk
$x\in\BbbR^r$. (i) Dengan menuliskan $\mu_r$ untuk ukuran Lebesgue pada
$\Bbb R^r$, tunjukkan bahwa
$\mu_r\phi^{-1}[E]=\int_Er\beta_rt^{r-1}\mu_1(dt)$
untuk setiap himpunan terukur Lebesgue $E\subseteq\coint{0,\infty}$. ({\it
Petunjuk\/}: mulailah dengan interval $E$, dengan memperhatikan dari 115Xe bahwa
$\mu_r\{x:\|x\|\le a\}=\beta_ra^r$ untuk $a\ge 0$, lalu lanjutkan ke himpunan
terbuka, himpunan terabaikan, dan himpunan terukur umum.) (ii) Dengan
menggunakan 235R, tunjukkan bahwa

$$\eqalign{\int e^{-\|x\|^2/2}\mu_r(dx)
&=r\beta_r\int_0^{\infty}t^{r-1}e^{-t^2/2}\mu_1(dt)
=2^{(r-2)/2}r\beta_r\Gamma(\Bover{r}2)\cr
&=2^{r/2}\beta_r\Gamma(1+\Bover{r}2)
=(\sqrt{2}\Gamma(\Bover12))^r\cr}$$

\noindent dengan $\Gamma$ merupakan fungsi-$\Gamma$ (225Xh). (iii) Tunjukkan
bahwa

\Centerline{$2\Gamma(\bover12)^2=2\beta_2\Gamma(2)
=2\beta_2\int_0^{\infty}te^{-t^2/2}dt=2\pi$,}

\noindent dan karena itu
$\beta_r=\Bover{\pi^{r/2}}{\Gamma(1+\bover{r}2)}$ dan
$\int_{-\infty}^{\infty}e^{-t^2/2}dt=\sqrt{2\pi}$.
%252Q

\leader{252Y}{Latihan lanjutan (a)}
%\spheader 252Ya
Misalkan $(X,\Sigma,\mu)$ ruang ukur. Tunjukkan bahwa
pernyataan-pernyataan berikut ekuivalen: (i) pelengkapan $\mu$ ditentukan
secara lokal; (ii) pelengkapan $\mu$ sama dengan versi c.l.d.\ dari
$\mu$; (iii) bilamana $(Y,\Tau,\nu)$ ruang ukur $\sigma$-hingga dan
$\lambda$ ukuran produk c.l.d.\ pada $X\times Y$, serta $f$ fungsi
sedemikian sehingga $\int fd\lambda$ terdefinisi dalam
$[-\infty,\infty]$, maka $\iint f(x,y)\nu(dy)\mu(dx)$ terdefinisi dan
sama dengan $\int fd\lambda$.
%252B

\spheader 252Yb Misalkan $(X,\Sigma,\mu)$ ruang ukur. Tunjukkan bahwa
pernyataan-pernyataan berikut ekuivalen: (i) $\mu$ memiliki himpunan-himpunan
terabaikan yang ditentukan secara lokal; (ii) bilamana $(Y,\Tau,\nu)$ ruang
ukur $\sigma$-hingga dan $\lambda$ ukuran produk c.l.d.\ pada
$X\times Y$, maka $\iint f(x,y)\nu(dy)\mu(dx)$ terdefinisi dan sama dengan
$\int fd\lambda$ untuk sebarang fungsi yang dapat diintegralkan terhadap
$\lambda$, katakanlah $f$.
%252Ya, 252B

\spheader 252Yc Misalkan $(X,\Sigma,\mu)$ dan $(Y,\Tau,\nu)$ adalah ruang
ukur, dan $\lambda_0$ ukuran produk primitif pada $X\times Y$ (251C).
Misalkan $f$ sebarang fungsi bernilai riil yang dapat diintegralkan terhadap
$\lambda_0$. Tunjukkan bahwa
$\iint f(x,y)\nu(dy)\mu(dx)=\int fd\lambda_0$. \Hint{tunjukkan bahwa terdapat
barisan-barisan $\sequencen{G_n}$, $\sequencen{H_n}$ dari himpunan-himpunan
berukuran hingga sedemikian sehingga $f(x,y)$ terdefinisi dan sama dengan $0$
untuk setiap
$(x,y)\in(X\times Y)\setminus\bigcup_{n\in\Bbb N}G_n\times H_n$.}
%252B

\spheader 252Yd Misalkan $(X,\Sigma,\mu)$ dan
$(Y,\Tau,\nu)$ adalah
ruang-ruang ukur; misalkan $\lambda_0$ ukuran produk primitif pada
$X\times Y$, dan $\lambda$ ukuran produk c.l.d. Tunjukkan bahwa jika
$f$ fungsi bernilai riil yang dapat diintegralkan terhadap $\lambda_0$, maka
fungsi tersebut dapat diintegralkan terhadap $\lambda$, dan
$\int fd\lambda=\int fd\lambda_0$.
%252B

\spheader 252Ye\dvAformerly{252Yj}
Misalkan $(X,\Sigma,\mu)$ ruang ukur yang lengkap dan ditentukan secara lokal,
serta $a<b$ dalam
$\Bbb R$, yang diperlengkapi dengan ukuran Lebesgue; misalkan $\Lambda$ domain
ukuran produk c.l.d.\ $\lambda$ pada $X\times[a,b]$. Misalkan
$f:X\times\ooint{a,b}\to\Bbb R$ fungsi terukur terhadap $\Lambda$ sedemikian
sehingga $t\mapsto f(x,t):[a,b]\to\Bbb R$ kontinu pada $[a,b]$ dan
terdiferensialkan pada $\ooint{a,b}$ untuk setiap $x\in X$. (i) Tunjukkan
bahwa turunan parsial $\Pd{f}{t}$ terhadap variabel kedua terukur terhadap
$\Lambda$. (ii) Sekarang andaikan bahwa $\Pd{f}{t}$ dapat diintegralkan
terhadap $\lambda$ dan bahwa $\int f(x,t_0)\mu(dx)$ terdefinisi dan berhingga
untuk suatu $t_0\in\ooint{a,b}$. Tunjukkan bahwa
$F(t)=\int f(x,t)\mu(dx)$ terdefinisi dalam $\Bbb R$ untuk setiap $t\in[a,b]$,
bahwa $F$ kontinu mutlak, dan bahwa
$F'(t)=\int\Pd{f}{t}(x,t)\mu(dx)$
untuk hampir setiap $t\in\ooint{a,b}$. ({\it Petunjuk\/}:
$F(c)=F(a)+\int_{X\times[a,c]}\Pd{f}{t}d\lambda$ untuk setiap $c\in[a,b]$.)
%252C %225L

\spheader 252Yf Tunjukkan bahwa
$\Bover{\Gamma(a)\Gamma(b)}{\Gamma(a+b)}=\int_0^1t^{a-1}(1-t)^{b-1}dt$
untuk semua $a$, $b>0$. ({\it Petunjuk\/}: tunjukkan bahwa

\Centerline{$\int_0^{\infty}t^{a-1}\int_t^{\infty}
  e^{-x}(x-t)^{b-1}dxdt
=\int_0^{\infty}e^{-x}\int_0^xt^{a-1}(x-t)^{b-1}dtdx$.)}
%252C

\spheader 252Yg Misalkan $(X,\Sigma,\mu)$ dan $(Y,\Tau,\nu)$ adalah
ruang-ruang ukur $\sigma$-hingga dan $\lambda$ ukuran produk c.l.d.\ pada
$X\times Y$. Andaikan bahwa
$f\in\eusm L^0(\lambda)$ dan bahwa $1<p<\infty$. Tunjukkan bahwa
$(\int|\int f(x,y)dx|^pdy)^{1/p}\le\int(\int|f(x,y)|^pdy)^{1/p}dx$.
({\it Petunjuk\/}:
tetapkan $q=\bover{p}{p-1}$ dan pertimbangkan integral
$\int|f(x,y)g(y)|\lambda(d(x,y))$ untuk $g\in\eusm L^q(\nu)$, dengan menggunakan
244K.)
%252C

\spheader 252Yh Misalkan $\nu$ ukuran Lebesgue pada $\coint{0,\infty}$;
andaikan bahwa $f\in\eusm L^p(\nu)$ dengan $1<p<\infty$. Tetapkan
$F(y)=\Bover1y\int_0^yf$ untuk $y>0$. Tunjukkan bahwa
$\|F\|_p\le\Bover{p}{p-1}\|f\|_p$. ({\it Petunjuk\/}:
$F(y)=\int_0^1f(xy)dx$; gunakan 252Yg dengan $X=[0,1]$,
$Y=\coint{0,\infty}$.)
%252Yg, 252C

\spheader 252Yi Tunjukkan bahwa jika $p$ adalah sebarang polinom (riil) tak nol
dalam $r$ variabel, maka $\{x:x\in\BbbR^r,\,p(x)=0\}$ terabaikan terhadap
ukuran Lebesgue.
%252D

\spheader 252Yj Misalkan $(X,\Sigma,\mu)$ dan $(Y,\Tau,\nu)$ adalah ruang
ukur, dengan produk c.l.d.\ $(X\times Y,\Lambda,\lambda)$. Misalkan $f$ fungsi
bernilai riil nonnegatif yang terukur terhadap $\Lambda$ dan didefinisikan
pada suatu himpunan koterabaikan terhadap $\lambda$, serta andaikan bahwa

\Centerline{$
\overline{\intop}\bigl(\overline{\intop}f(x,y)\mu(dx)\bigr)
\nu(dy)$}

\noindent berhingga. Tunjukkan bahwa $f$ dapat diintegralkan terhadap
$\lambda$.
%252G

\spheader 252Yk Misalkan $(X,\Sigma,\mu)$ adalah interval satuan $[0,1]$
dengan ukuran Lebesgue, dan $(Y,\Tau,\nu)$ interval tersebut dengan ukuran
cacah, sebagaimana dalam 252K; misalkan $\lambda_0$ ukuran produk primitif pada
$[0,1]^2$. (i) Dengan menetapkan $\Delta=\{(t,t):t\in[0,1]\}$, tunjukkan
bahwa $\lambda_0\Delta=\infty$. (ii) Tunjukkan bahwa
$\lambda_0$ tidak semihingga. (iii) Tunjukkan bahwa jika
$W\in\dom\lambda_0$, maka
$\lambda_0W=\sum_{y\in[0,1]}\mu W^{-1}[\{y\}]$
jika terdapat himpunan terhitung $A\subseteq[0,1]$ dan himpunan terabaikan
Lebesgue $E\subseteq[0,1]$ sedemikian sehingga
$W\subseteq([0,1]\times A)\cup(E\times[0,1])$, dan bernilai $\infty$ jika tidak.
%252K

\spheader 252Yl Misalkan $(X,\Sigma,\mu)$ ruang ukur, dan
$\lambda_0$ ukuran produk primitif pada
$X\times\Bbb R$, dengan $\Bbb R$ diberi ukuran Lebesgue; tuliskan
$\Lambda$ untuk domainnya. Untuk sebarang fungsi bernilai
$[0,\infty]$, katakanlah $f$, yang didefinisikan pada subhimpunan koterabaikan dari
$X$, tuliskan $\Omega_f$, $\Omega'_f$ untuk himpunan-himpunan ordinat yang
bersesuaian, sebagaimana dalam 252N. Tunjukkan bahwa jika salah satu dari
$\lambda_0\Omega_f$, $\lambda_0\Omega'_f$,
$\int fd\mu$ terdefinisi dan berhingga, demikian pula yang lainnya, dan
ketiganya sama.
%252B, 252N

\spheader 252Ym Misalkan $(X,\Sigma,\mu)$ ruang ukur yang lengkap dan
ditentukan secara lokal, dan $f$ fungsi nonnegatif yang didefinisikan pada
subhimpunan koterabaikan dari $X$. Tuliskan $\Omega_f$, $\Omega'_f$ untuk
himpunan-himpunan ordinat yang bersesuaian, sebagaimana dalam 252N. Misalkan
$\lambda$ ukuran produk c.l.d.\ pada $X\times\Bbb R$, dengan
$\Bbb R$ diberi ukuran Lebesgue. Tunjukkan bahwa
$\overline{\int}f\,d\mu=\lambda^*\Omega_f=\lambda^*\Omega'_f$.
%252N, 252B

\spheader 252Yn Misalkan $(X,\Sigma,\mu)$ ruang ukur dan
$f:X\to\coint{0,\infty}$ fungsi. Tunjukkan bahwa
$\overline{\int}fd\mu=\int_0^{\infty}\mu^*\{x:f(x)\ge t\}dt$.
%252O

\spheader 252Yo Misalkan $(X,\Sigma,\mu)$ ruang ukur lengkap dan
tuliskan $\eusm M^{0,\infty}$ untuk himpunan
$\{f:f\in\eusm L^0(\mu),\,\mu\{x:|f(x)|\ge a\}$ berhingga untuk suatu
$a\in\coint{0,\infty}\}$. (i) Tunjukkan bahwa untuk setiap
$f\in\eusm M^{0,\infty}$
terdapat fungsi tak-naik $f^*:\ooint{0,\infty}\to\Bbb R$ sedemikian sehingga
$\mu_L\{t:f^*(t)\ge\alpha\}=\mu\{x:|f(x)|\ge\alpha\}$ untuk
setiap $\alpha>0$, dengan menuliskan $\mu_L$ untuk ukuran Lebesgue. (ii) Tunjukkan
bahwa $\int_E|f|d\mu\le\int_0^{\mu E}f^*d\mu_L$ untuk setiap $E\in\Sigma$
(dengan mengizinkan $\infty$). ({\it Petunjuk\/}: $(f\times\chi E)^*\le f^*$.)
(iii) Tunjukkan bahwa $\|f^*\|_p=\|f\|_p$ untuk setiap $p\in[1,\infty]$,
$f\in\eusm M^{0,\infty}$. ({\it
Petunjuk\/}: $(|f|^p)^*=(f^*)^p$.) (iv) Tunjukkan bahwa jika $f$,
$g\in\eusm M^{0,\infty}$ maka
$\int|f\times g|d\mu\le\int f^*\times g^*d\mu_L$. ({\it
Petunjuk\/}: tinjau fungsi-fungsi sederhana terlebih dahulu.) (v) Tunjukkan bahwa
jika $\mu$ tanpa atom maka
$\int_0^af^*d\mu_L=\sup_{E\in\Sigma,\mu E\le a}\int_E|f|$
untuk setiap $a\ge 0$. ({\it Petunjuk\/}: 215D.) (vi) Tunjukkan bahwa
$A\subseteq\eusm L^1(\mu)$ terintegralkan secara seragam jika dan hanya jika
$\{f^*:f\in A\}$ terintegralkan secara seragam dalam $\eusm L^1(\mu_L)$.
($f^*$ disebut {\bf penyusunan ulang menurun} dari $f$.)
%252O

\spheader 252Yp Misalkan $(X,\Sigma,\mu)$ ruang ukur yang lengkap dan
ditentukan secara lokal, dan tuliskan $\nu$ untuk ukuran Lebesgue pada
$[0,1]$. Tunjukkan bahwa ukuran produk c.l.d.\ $\lambda$ pada
$X\times[0,1]$ dapat dilokalkan jika dan hanya jika $\mu$ dapat dilokalkan.
({\it Petunjuk-petunjuk\/}: (i) jika
$\Cal E\subseteq\Sigma$, tunjukkan bahwa $F\in\Sigma$ merupakan supremum
esensial untuk $\Cal E$ dalam $\Sigma$ jika dan hanya jika $F\times[0,1]$
merupakan supremum esensial untuk
$\{E\times[0,1]:E\in\Cal E\}$ dalam $\Lambda=\dom\lambda$. (ii) Untuk
$W\in\Lambda$, $n\in\Bbb N$, $k<2^n$ tetapkan

\Centerline{$W_{nk}
=\{x:x\in X,\,\nu^*\{t:(x,t)\in W,\,2^{-n}k\le t\le 2^{-n}(k+1)\}
\ge 2^{-n-1}\}$.}

\noindent Tunjukkan bahwa jika $\Cal W\subseteq\Lambda$ dan $F_{nk}$ merupakan
supremum esensial untuk $\{W_{nk}:W\in\Cal W\}$ dalam $\Sigma$ untuk semua
$n$, $k$, maka

\Centerline{$\bigcup_{n\in\Bbb N}\bigcap_{m\ge n}
\bigcup_{k<2^{m}}F_{mk}\times[2^{-m}k,2^{-m}(k+1)]$}

\noindent merupakan supremum esensial untuk $\Cal W$ dalam $\Lambda$.)
%252O

\spheader 252Yq Misalkan $(X,\Sigma,\mu)$ adalah ruang dari Contoh
216D, dan berikan ukuran Lebesgue pada $[0,1]$. Tunjukkan bahwa ukuran produk
c.l.d.\ pada $X\times [0,1]$ lengkap, ditentukan secara lokal,
tanpa atom, dan tidak dapat dilokalkan.
%252Yp 252O

\spheader 252Yr Misalkan $(X,\Sigma,\mu)$ ruang ukur yang lengkap dan
ditentukan secara lokal dan $(Y,\Tau,\nu)$ ruang ukur semihingga dengan
$\nu Y>0$. Tunjukkan bahwa jika ukuran produk c.l.d.\ pada
$X\times Y$ dapat dilokalkan secara ketat, maka $\mu$ dapat dilokalkan secara
ketat. ({\it Petunjuk\/}: ambil $F\in\Tau$, $0<\nu F<\infty$. Misalkan
$\langle W_i\rangle_{i\in I}$ dekomposisi dari $X\times Y$. Untuk
$i\in I$, $n\in\Bbb N$
tetapkan $E_{in}=\{x:\nu^*\{y:y\in F,\,(x,y)\in W_i\}\ge 2^{-n}\}$. Terapkan
213Yf pada $\{E_{in}:i\in I,\,n\in\Bbb N\}$.)
%252Yq 252O

\spheader 252Ys Misalkan $(X,\Sigma,\mu)$ adalah ruang dari Contoh
216E, dan berikan ukuran Lebesgue pada $[0,1]$. Tunjukkan bahwa ukuran produk
c.l.d.\ pada $X\times[0,1]$ lengkap, ditentukan secara lokal,
tanpa atom dan dapat dilokalkan, tetapi tidak dapat dilokalkan secara ketat.
%252Yr, 252O

\spheader 252Yt Misalkan $(X,\Sigma,\mu)$ ruang ukur dan $f$ fungsi bernilai
kompleks yang dapat diintegralkan terhadap $\mu$. Untuk
$\alpha\in\ocint{-\pi,\pi}$
tetapkan $H_{\alpha}=\{x:x\in\dom f,\,\Real(e^{-i\alpha}f(x))>0\}$. Tunjukkan
bahwa $\int_{-\pi}^{\pi}\Real(e^{-i\alpha}\int_{H_{\alpha}}f)d\alpha
\discrversionA{\break}{}=2\int|f|$, dan karena itu terdapat suatu
$\alpha$ sedemikian sehingga
$|\int_{H_{\alpha}}f|\ge\Bover1{\pi}\int|f|$. (Bandingkan 246F.)
%246F, 252R

\spheader 252Yu Tetapkan $f(t)=t-\ln(t+1)$ untuk $t>-1$. (i) Tunjukkan bahwa
$\Gamma(a+1)=a^{a+1}e^{-a}\int_{-1}^{\infty}e^{-af(u)}du$ untuk setiap
$a>0$. \Hint{substitusikan $u=\bover{t}{a}-1$ dalam 225Xh(iii).} (ii) Tunjukkan
bahwa terdapat $\delta>0$ sedemikian sehingga $f(t)\ge\bover13t^2$ untuk
$-1\le t\le\delta$. (iii)
Dengan menetapkan $\alpha=\bover12f(\delta)$, tunjukkan bahwa (untuk $a\ge 1$)

\Centerline{$\sqrt{a}\int_{\delta}^{\infty}e^{-af(t)}dt
\le\sqrt{a}e^{-a\alpha}\int_0^{\infty}e^{-f(t)/2}dt\to 0$}

\noindent ketika $a\to\infty$. (iv) Tetapkan
$g_a(t)=e^{-af(t/\sqrt{a})}$ jika
$-\sqrt{a}<t\le\delta\sqrt{a}$, dan $0$ jika tidak. Tunjukkan bahwa
$g_a(t)\le e^{-t^2/3}$ untuk semua $a$, $t$ dan bahwa
$\lim_{a\to\infty}g_a(t)=e^{-t^2/2}$ untuk semua $t$, sehingga

$$\eqalign{\lim_{a\to\infty}\Bover{e^a\Gamma(a+1)}{a^{a+\bover12}}
&=\lim_{a\to\infty}\sqrt{a}\int_{-1}^{\infty}e^{-af(t)}dt
=\lim_{a\to\infty}\sqrt{a}\int_{-1}^{\delta}e^{-af(t)}dt\cr
&=\lim_{a\to\infty}\int_{-\infty}^{\infty}g_a(t)dt
=\int_{-\infty}^{\infty}e^{-t^2/2}dt=\sqrt{2\pi}.\cr}$$

\noindent (v) Tunjukkan bahwa
$\lim_{n\to\infty}\Bover{n!}{e^{-n}n^n\sqrt{n}}=\sqrt{2\pi}$.
(Ini adalah {\bf rumus Stirling}.)
%252+

\spheader 252Yv Misalkan $(X,\Sigma,\mu)$ ruang ukur yang lengkap dan
ditentukan secara lokal, serta $f$, $g$ dua fungsi bernilai riil yang terukur
secara virtual terhadap $\mu$ dan didefinisikan hampir di mana-mana dalam $X$.
(i) Misalkan $\lambda$ produk c.l.d.\ dari $\mu$ dan ukuran Lebesgue pada
$\Bbb R$.
Dengan menetapkan $\Omega^*_f=\{(x,a):x\in\dom f$, $a\in\Bbb R$, $a\le f(x)\}$ dan
$\Omega^*_g=\{(x,a):x\in\dom g$, $a\in\Bbb R$, $a\le g(x)\}$, tunjukkan bahwa
$\lambda(\Omega^*_f\setminus\Omega^*_g)=\int(f-g)^+d\mu$ dan
$\lambda(\Omega^*_f\symmdiff\Omega^*_g)=\int|f-g|d\mu$.
(ii) Andaikan bahwa $\mu$ $\sigma$-hingga. Tunjukkan bahwa

\Centerline{$\int|f-g|d\mu
=\int_{-\infty}^{\infty}
\mu(\{x:x\in\dom f\cap\dom g,\,
  (f(x)-a)(g(x)-a)<0\})\,da$.}

\noindent(iii) Andaikan bahwa $\mu$ $\sigma$-hingga, bahwa $\Tau$ merupakan
subaljabar-sigma dari $\Sigma$, bahwa $E\in\Sigma$ dan bahwa $g:X\to[0,1]$
terukur terhadap $\Tau$. Tunjukkan bahwa terdapat $F\in\Tau$ sedemikian sehingga
$\mu(E\symmdiff F)\le\int|\chi E-g|d\mu$.
%252N
}%end of exercises

\endnotes{
\Notesheader{252} Selama satu setengah jilid ini saya telah meminta Anda
menerima gagasan mengintegralkan fungsi yang terdefinisi sebagian, dengan
bersikeras bahwa cepat atau lambat fungsi-fungsi tersebut akan muncul di inti
pokok bahasan ini. Saat itu kini telah tiba. Jika kita hendak menerapkan teorema
Fubini dan Tonelli dalam kasus yang paling mendasar, dengan kedua faktor sama
dengan ukuran Lebesgue pada interval satuan, tentulah wajar untuk meninjau
semua fungsi yang dapat diintegralkan pada persegi terhadap ukuran Lebesgue
berdimensi dua. Ukuran Lebesgue berdimensi dua adalah ukuran lengkap, sehingga,
khususnya, memberikan ukuran nol kepada setiap himpunan berbentuk
$\{(x,b):x\in A\}$ atau $\{(a,y):y\in A\}$, terlepas dari apakah himpunan
$A$ diukur oleh ukuran berdimensi satu. Oleh sebab itu, jika $f$ merupakan
fungsi dua variabel yang dapat diintegralkan terhadap ukuran Lebesgue
berdimensi dua, tidak ada alasan mengapa sebarang penampang tertentu
$x\mapsto f(x,b)$ atau $y\mapsto f(a,y)$ harus terukur, apalagi dapat
diintegralkan. Akibatnya, sekalipun $f$ sendiri didefinisikan di mana-mana,
integral luar dari $\iint f(x,y)dxdy$ mungkin diterapkan kepada suatu fungsi
yang tidak didefinisikan untuk setiap $y$.
Perlu saya tegaskan bahwa persoalannya bukanlah `$\infty$'; fungsi-fungsi
yang menyulitkan ialah yang penampangnya begitu tak teratur sehingga integral
sama sekali tidak dapat diberikan kepadanya.

Saya telah melihat banyak pendekatan terhadap persoalan pelik khusus ini,
yang pada umumnya kurang sepenuh hati daripada pendekatan yang saya tetapkan
untuk risalah ini. Sebagian kesulitannya ialah bahwa teorema Fubini sungguh
berada di pusat teori ukuran. Pada bagian-bagian luas dari pokok bahasan ini,
dapat dikatakan bahwa suatu hasil tidak sepele jika dan hanya jika hasil itu
bergantung pada teorema Fubini. Karena itu saya enggan memasukkan tambalan
lokal apa pun dengan mengatakan bahwa `dalam bab ini, kita akan
mengintegralkan fungsi yang tidak didefinisikan di mana-mana'; tak lama
kemudian, ketentuan semacam itu harus disisipkan ke dalam pembukaan separuh
teorema terbaik, atau harus diberikan penjelasan mengapa ketentuan itu tidak
diperlukan dalam konteks masing-masing. Saya kira salah satu tanggapan yang
paling umum ialah (seperti {\smc Halmos 50}) membatasi perhatian pada
fungsi yang terukur terhadap $\Sigma\tensorhat\Tau$, yang untuk sementara
menghilangkan masalah keterukuran (252Xh, 252P); tetapi sayangnya (atau,
menurut saya, justru untungnya) terdapat penerapan-penerapan penting yang
fungsinya sebenarnya tidak terukur terhadap $\Sigma\tensorhat\Tau$, melainkan
termasuk dalam suatu kelas yang lebih luas, dan cepat atau lambat pembatasan
ini menimbulkan akrobat yang tidak elok ketika kita terpaksa menyesuaikan hasil
terbatas dengan konteks yang tidak terduga. Selain itu, pembatasan tersebut
tidak menyatakan informasi yang sesungguhnya cukup penting, yakni bahwa jika
$f$ fungsi terukur dari dua variabel, maka (di bawah syarat-syarat yang sesuai)
hampir semua penampangnya terukur (252E).

Dalam 252B dan akibat-akibatnya terdapat pembatasan yang canggung: kita
mengandaikan bahwa salah satu ukuran $\sigma$-hingga dan ukuran yang lain dapat
dilokalkan secara ketat atau lengkap dan ditentukan secara lokal. Pertanyaan
yang jelas ialah apakah hipotesis-hipotesis ini diperlukan. Dari 252K kita
melihat bahwa hipotesis `$\sigma$-hingga' pada faktor kedua tentu tidak dapat
ditinggalkan, bahkan ketika faktor pertama merupakan ukuran probabilitas
lengkap. Persyaratan `$\mu$ dapat dilokalkan secara ketat atau lengkap dan
ditentukan secara lokal' sebenarnya sedikit lebih kuat daripada yang
diperlukan, dan juga terlampau rumit. Hipotesis yang `tepat' ialah bahwa
pelengkapan $\mu$ harus ditentukan secara lokal (lihat 252Ya). Intinya ialah
bahwa, karena produk dua ukuran sama dengan produk versi c.l.d.\ masing-masing
(251T), tidak ada teorema yang beranjak dari ukuran produk ke ukuran faktor
yang dapat membedakan suatu ukuran dari versi c.l.d.-nya; sehingga, dalam
252B, kita harus menduga bahwa $\mu$ dan versi c.l.d.-nya perlu menghasilkan
integral yang sama. Pembuktian 252B akan lebih terarah jika hipotesisnya
disederhanakan menjadi `$\nu$ $\sigma$-hingga dan $\mu$ lengkap serta
ditentukan secara lokal'. Tetapi ini hanya akan memindahkan sebagian argumen
ke dalam pembuktian 252C.

Kita juga harus bekerja sedikit lebih keras dalam 252B agar mencakup fungsi
dan integral yang mengambil nilai $\pm\infty$. Teorema Fubini begitu sentral
bagi teori ukuran sehingga saya yakin sedikit upaya tambahan layak dilakukan
untuk menyatakan hasil-hasilnya dengan keumuman maksimum. Hal ini terutama
penting karena kita sering menerapkannya pada integral berulang majemuk, seperti
dalam 252Xd, ketika kendali kita terhadap fungsi-fungsi antara yang harus
diintegralkan bahkan lebih kecil daripada biasanya.

Saya telah mengungkapkan semua hasil utama bagian ini dalam istilah ukuran
produk `c.l.d.'. Dalam kasus ruang-ruang $\sigma$-hingga, tentu saja, tempat
teorinya bekerja paling baik, kita dapat pula menggunakan ukuran produk
`primitif'. Bahkan teorema Fubini sendiri mempunyai versi dalam istilah
ukuran produk primitif yang agak lebih anggun daripada 252B sebagaimana
dinyatakan (252Yc), dan mencakup sebagian besar penerapan. (Integral terhadap
ukuran produk primitif dan c.l.d.\ tentu saja sangat erat hubungannya; lihat
252Yd.) Namun kadang-kadang kita memang perlu meninjau ruang yang tidak
$\sigma$-hingga, dan dalam kasus-kasus ini bentuk tak simetris dalam 252B
mendekati bentuk terbaik yang dapat kita peroleh. Menggunakan ukuran produk
primitif sama sekali tidak membantu mengatasi hambatan paling besar, yaitu
fenomena dalam 252K (lihat 252Yk).

Konsep prakalkulus tentang integral sebagai `luas di bawah kurva' diwujudkan
dalam 252N: integral fungsi nonnegatif merupakan ukuran himpunan ordinatnya.
Ini tidak memuaskan sebagai definisi integral, bukan hanya karena persyaratan
bahwa ruang dasar harus lengkap dan ditentukan secara lokal (yang dapat
ditangani dengan menggunakan ukuran produk primitif, seperti dalam 252Yl),
melainkan juga karena konstruksi ukuran produk melibatkan integrasi (bagian
(c) dari pembuktian 251E).
Gagasan 252N ialah menghubungkan ukuran suatu himpunan ordinat dengan integral
ukuran penampang-penampang vertikalnya. Anehnya, jika sebagai gantinya kita
mengintegralkan ukuran penampang-penampang {\it horizontal}, seperti dalam
252O, kita memperoleh hasil yang lebih luwes. (Memang hasil ini tidak
melibatkan konsep `ukuran produk', dan dapat muncul pada sebarang titik setelah
\S123.) Perhatikan bahwa integral $\int_0^{\infty}\ldots dt$ di sini
diterapkan kepada fungsi monoton, sehingga dapat ditafsirkan sebagai integral
Riemann tak wajar. Jika Anda merasa cukup memahami integral Riemann sehingga
hal ini menjadi alternatif yang menggoda bagi konstruksi dalam \S122, bagian
yang sulit sekarang ialah pembuktian bahwa integral tersebut aditif.

Jalur argumen yang berbeda ialah menggunakan integrasi atas penampang untuk
mendefinisikan ukuran produk. Kesulitan pendekatan ini ialah bahwa, kecuali
kita sangat berhati-hati, kita mungkin mendapati konstruksi yang tak simetris.
Menurut pandangan saya, ketaksimetrian semacam itu dapat diterima hanya jika
tidak ada alternatif. Namun dalam Bab 43 Jilid 4 saya akan mendeskripsikan
sepasang contoh.

Dari dua contoh yang saya berikan di sini, 252K dimaksudkan untuk menunjukkan
bahwa ketika saya mensyaratkan ruang-ruang $\sigma$-hingga, ruang-ruang itu
memang benar-benar diperlukan, sedangkan 252L dimaksudkan untuk menunjukkan
bahwa keterukuran bersama bersifat esensial dalam teorema Tonelli dan
akibat-akibatnya. Ruang-ruang faktor dalam 252K, yaitu ukuran Lebesgue dan
ukuran cacah, dipilih untuk menunjukkan bahwa hanya ketiadaan sifat
$\sigma$-hingga yang dapat menjadi masalah; selain itu, keduanya seteratur
yang secara wajar dapat diminta. Dalam 252L saya menggunakan ukuran
terhitung-koterhitung pada $\omega_1$, yang mungkin Anda anggap hanya pantas
untuk contoh tandingan; dan memang timbul pertanyaan apakah fenomena yang sama
terjadi pada ukuran Lebesgue. Ini membawa kita ke persoalan yang mendalam, dan
saya akan kembali kepadanya dalam Bab 53 Jilid 5.

Barangkali perlu saya catat secara eksplisit bahwa dalam teorema Fubini, kita
benar-benar memerlukan fungsi yang dapat diintegralkan terhadap ukuran produk.
Saya menyertakan 252Xf dan 252Xg untuk mengingatkan Anda bahwa bahkan dalam
keadaan yang paling teratur, integral berulang $\iint f\,dxdy$,
$\iint f\,dydx$ dapat tidak sama jika $f$ tidak dapat diintegralkan sebagai
fungsi dua variabel.

Ada banyak cara untuk menghitung volume $\beta_r$ dari bola berdimensi $r$;
cara yang saya gunakan dalam 252Q mengikuti jalur yang sudah terasa alami bagi
saya sebelum pernah mendengar teori ukuran. Dalam 252Xi saya mengusulkan
metode lain. Gagasan integrasi dengan substitusi, yang digunakan dalam bagian
(b) dari argumen untuk 252Q, di sana ditopang oleh argumen ad hoc; saya akan
menyajikan pendekatan lain yang dapat diterapkan lebih umum dalam Bab 26. Di
tempat lain (252Xi, 252Yf, 252Yh, 252Yu) saya mendapati diri saya menerima
begitu saja substitusi berbentuk $t\mapsto at$, $t\mapsto a+t$,
$t\mapsto t^2$; untuk pembenaran sistematis, lihat \S263. Tentu saja sangat
banyak rumus kalkulus lanjut lainnya juga didasarkan pada integrasi berulang
dari satu jenis atau jenis lainnya, dan saya memberikan segelintir contoh
hasil semacam itu (252Xb,
252Ye-252Yh, %252Ye 252Yf 252Yg 252Yh
252Yu).

}%end of notes

\discrpage

\leaveitout{%\spheader 252Y?
Misalkan $(X,\Sigma,\mu)$ dan $(Y,\Tau,\nu)$ adalah ruang-ruang ukur, dan
$\lambda_0$
ukuran produk primitif pada $X\times Y$. Misalkan $f$ fungsi bernilai riil
yang didefinisikan $\lambda_0$-hampir di mana-mana pada $X\times Y$. Tunjukkan
bahwa

\Centerline{$
\overline{\intop}\bigl(\overline{\intop}f(x,y)\mu(dx)\bigr)\nu(dy)
\le\overline{\intop}f\,d\lambda_0$.}
%252B
}%end of leaveitout
